\documentclass[UTF8,11pt,a4paper]{ctexart}
\usepackage[margin=2.4cm,top=2.7cm,bottom=2.5cm]{geometry}
\usepackage{enumitem}
\usepackage{xcolor}
\usepackage{hyperref}
\usepackage{fancyhdr}
\usepackage{setspace}
\usepackage{titlesec}
\usepackage{ulem}
\usepackage{CJKfntef}
\definecolor{linkblue}{HTML}{1A48A8}
\hypersetup{colorlinks=true,linkcolor=linkblue,urlcolor=linkblue}
\setstretch{1.45}
\setlength{\parindent}{2em}
\setlength{\parskip}{0.25em}
\titleformat{\section}{\centering\Large\bfseries}{\thesection}{0.6em}{}
\titleformat{\subsection}{\large\bfseries}{\thesubsection}{0.5em}{}
\pagestyle{fancy}
\fancyhf{}
\fancyhead[C]{\small 入党积极分子培训班基础知识学习资料}
\fancyfoot[C]{\small\thepage}
\renewcommand{\headrulewidth}{0.4pt}
\setlist[enumerate]{leftmargin=2.3em,itemsep=0.25em,topsep=0.3em}
\setlist[itemize]{leftmargin=2em,itemsep=0.15em,topsep=0.2em}
\newcommand{\blank}{\underline{\hspace{4em}}}
\begin{document}
\begin{titlepage}
\centering
\vspace*{5cm}
{\Huge\bfseries 入党积极分子培训班}\\[0.8em]
{\Huge\bfseries 基础知识学习资料}\\[2.5em]
{\large 北航学生入党积极分子培训班}\\[1em]
{\large 2026 年 10 月}
\vfill
\end{titlepage}

\section*{前言}

本资料为党的相关基础知识资料，大家可结合此资料内容，认真学习《高举中国特色社会主义伟大旗帜，为全面建设社会主义现代化国家而团结奋斗》、《中国共产党章程》（中国共产党第二十次全国代表大会部分修改，2022年10月）和《中国共产党发展党员工作细则》（2026年5月），认真学习马克思主义基本理论、党的基本理论、党史、党纪等。建议关注学习强国平台中马克思主义理论、党建、党史、党纪等方面相关内容。望大家常学常练、温故知新，力争学深悟透，不断提升理论认知和素养。

\section*{第一部分}

\subsection*{一、党的二十大报告（节选示例）}

\begin{enumerate}[label=\arabic*., leftmargin=2.3em]
\item 报告共分为15个部分：
\begin{itemize}[leftmargin=2.2em]
\item 过去五年的工作和新时代十年的伟大变革；
\item 开辟马克思主义中国化时代化新境界；
\item 新时代新征程中国共产党的使命任务；
\item 加快构建新发展格局，着力推动高质量发展；
\item 实施科教兴国战略，强化现代化建设人才支撑；
\item 发展全过程人民民主，保障人民当家作主；
\item 坚持全面依法治国，推进法治中国建设；
\item 推进文化自信自强，铸就社会主义文化新辉煌；
\item 增进民生福祉，提高人民生活品质；
\item 推动绿色发展，促进人与自然和谐共生；
\item 推进国家安全体系和能力现代化，坚决维护国家安全和社会稳定；
\item 实现建军一百年奋斗目标，开创国防和军队现代化新局面；
\item 坚持和完善“一国两制”，推进祖国统一；
\item 促进世界和平与发展，推动构建人类命运共同体；
\item 坚定不移全面从严治党，深入推进新时代党的建设新的伟大工程。
\end{itemize}
\item 马克思主义是我们立党立国、兴党兴国的根本指导思想。实践告诉我们，中国共产党为什么能，中国特色社会主义为什么好，归根到底是马克思主义行，是中国化时代化的马克思主义行。
\item 中国共产党勇于进行理论探索和创新，以全新的视野深化对共产党执政规律、社会主义建设规律、人类社会发展规律的认识，取得重大理论创新成果，集中体现为习近平新时代中国特色社会主义思想。
\item 教育、科技、人才是全面建设社会主义现代化国家的基础性、战略性支撑。必须坚持科技是第一生产力、人才是第一资源、创新是第一动力，深入实施科教兴国战略、人才强国战略、创新驱动发展战略，开辟发展新领域新赛道，不断塑造发展新动能新优势。
\item 我国是工人阶级领导的、以工农联盟为基础的人民民主专政的社会主义国家，国家一切权力属于人民。人民民主是社会主义的生命，是全面建设社会主义现代化国家的应有之义。全过程人民民主是社会主义民主政治的本质属性，是最广泛、最真实、最管用的民主。必须坚定不移走中国特色社会主义政治发展道路，坚持党的领导、人民当家作主、依法治国有机统一，坚持人民主体地位，充分体现人民意志、保障人民权益、激发人民创造活力。
\item 全面依法治国是国家治理的一场深刻革命，关系党执政兴国，关系人民幸福安康，关系党和国家长治久安。必须更好发挥法治固根本、稳预期、利长远的保障作用，在法治轨道上全面建设社会主义现代化国家。
\item 全面建设社会主义现代化国家，必须坚持中国特色社会主义文化发展道路，增强文化自信，围绕举旗帜、聚民心、育新人、兴文化、展形象建设社会主义文化强国，发展面向现代化、面向世界、面向未来的，民族的科学的大众的社会主义文化，激发全民族文化创新创造活力，增强实现中华民族伟大复兴的精神力量。
\item 江山就是人民，人民就是江山。中国共产党领导人民打江山、守江山，守的是人民的心。治国有常，利民为本。为民造福是立党为公、执政为民的本质要求。必须坚持在发展中保障和改善民生，鼓励共同奋斗创造美好生活，不断实现人民对美好生活的向往。
\item 尊重自然、顺应自然、保护自然，是全面建设社会主义现代化国家的内在要求。必须牢固树立和践行绿水青山就是金山银山的理念，站在人与自然和谐共生的高度谋划发展。
\item 国家安全是民族复兴的根基，社会稳定是国家强盛的前提。必须坚定不移贯彻总体国家安全观，把维护国家安全贯穿党和国家工作各方面全过程，确保国家安全和社会稳定。
\item 如期实现建军一百年奋斗目标，加快把人民军队建成世界一流军队，是全面建设社会主义现代化国家的战略要求。
\item “一国两制”是中国特色社会主义的伟大创举，是香港、澳门回归后保持长期繁荣稳定的最佳制度安排，必须长期坚持。
\item 中国始终坚持维护世界和平、促进共同发展的外交政策宗旨，致力于推动构建人类命运共同体。中国坚定奉行独立自主的和平外交政策，始终根据事情本身的是非曲直决定自己的立场和政策，维护国际关系基本准则，维护国际公平正义。
\item 全面建设社会主义现代化国家、全面推进中华民族伟大复兴，关键在党。
\item 全党必须牢记，全面从严治党永远在路上，党的自我革命永远在路上，决不能有松劲歇脚、疲劳厌战的情绪，必须持之以恒推进全面从严治党，深入推进新时代党的建设新的伟大工程，以党的自我革命引领社会革命。
\end{enumerate}

\subsection*{二、中国共产党章程（节选示例）}

\begin{enumerate}[label=\arabic*., leftmargin=2.3em]
\item 中国共产党是中国工人阶级的先锋队，同时是中国人民和中华民族的先锋队，是中国特色社会主义事业的领导核心，代表中国先进生产力的发展要求，代表中国先进文化的前进方向，代表中国最广大人民的根本利益。党的最高理想和最终目标是实现共产主义。
\item 中国共产党以马克思列宁主义、毛泽东思想、邓小平理论、“三个代表”重要思想、科学发展观、习近平新时代中国特色社会主义思想作为自己的行动指南。
\item 党是根据自己的纲领和章程，按照民主集中制组织起来的统一整体。
\item 党的最高领导机关，是党的全国代表大会和它所产生的中央委员会。党的地方各级领导机关，是党的地方各级代表大会和它们所产生的委员会。党的各级委员会向同级的代表大会负责并报告工作。
\item 党的全国代表大会的职权是：
\begin{itemize}[leftmargin=2.2em,label={}]
\item （一）听取和审查中央委员会的报告；
\item （二）审查中央纪律检查委员会的报告；
\item （三）讨论并决定党的重大问题；
\item （四）修改党的章程；
\item （五）选举中央委员会；
\item （六）选举中央纪律检查委员会。
\end{itemize}
\item 以毛泽东同志为主要代表的中国共产党人，把马克思列宁主义的基本原理同中国革命的具体实践结合起来，创立了毛泽东思想。毛泽东思想是马克思列宁主义在中国的运用和发展，是被实践证明了的关于中国革命和建设的正确的理论原则和经验总结，是中国共产党集体智慧的结晶。在毛泽东思想指引下，中国共产党领导全国各族人民，经过长期的反对帝国主义、封建主义、官僚资本主义的革命斗争，取得了新民主主义革命的胜利，建立了人民民主专政的中华人民共和国；新中国成立以后，顺利地进行了社会主义改造，完成了从新民主主义到社会主义的过渡，确立了社会主义基本制度，发展了社会主义的经济、政治和文化。
\item 十一届三中全会以来，以邓小平同志为主要代表的中国共产党人，总结新中国成立以来正反两方面的经验，解放思想，实事求是，实现全党工作中心向经济建设的转移，实行改革开放，开辟了社会主义事业发展的新时期，逐步形成了建设中国特色社会主义的路线、方针、政策，阐明了在中国建设社会主义、巩固和发展社会主义的基本问题，创立了邓小平理论。邓小平理论是马克思列宁主义的基本原理同当代中国实践和时代特征相结合的产物，是毛泽东思想在新的历史条件下的继承和发展，是马克思主义在中国发展的新阶段，是当代中国的马克思主义，是中国共产党集体智慧的结晶，引导着我国社会主义现代化事业不断前进。
\item 十三届四中全会以来，以江泽民同志为主要代表的中国共产党人，在建设中国特色社会主义的实践中，加深了对什么是社会主义、怎样建设社会主义和建设什么样的党、怎样建设党的认识，积累了治党治国新的宝贵经验，形成了“三个代表”重要思想。“三个代表”重要思想是对马克思列宁主义、毛泽东思想、邓小平理论的继承和发展，反映了当代世界和中国的发展变化对党和国家工作的新要求，是加强和改进党的建设、推进我国社会主义自我完善和发展的强大理论武器，是中国共产党集体智慧的结晶，是党必须长期坚持的指导思想。始终做到“三个代表”，是中国共产党的立党之本、执政之基、力量之源。
\item 十六大以来，以胡锦涛同志为主要代表的中国共产党人，坚持以邓小平理论和“三个代表”重要思想为指导，根据新的发展要求，深刻认识和回答了新形势下实现什么样的发展、怎样发展等重大问题，形成了以人为本、全面协调可持续发展的科学发展观。科学发展观是同马克思列宁主义、毛泽东思想、邓小平理论、“三个代表”重要思想既一脉相承又与时俱进的科学理论，是马克思主义关于发展的世界观和方法论的集中体现，是马克思主义中国化重大成果，是中国共产党集体智慧的结晶，是发展中国特色社会主义必须长期坚持的指导思想。
\item 十八大以来，以习近平同志为主要代表的中国共产党人，坚持把马克思主义基本原理同中国具体实际相结合、同中华优秀传统文化相结合，科学回答了新时代坚持和发展什么样的中国特色社会主义、怎样坚持和发展中国特色社会主义等重大时代课题，创立了习近平新时代中国特色社会主义思想。习近平新时代中国特色社会主义思想是对马克思列宁主义、毛泽东思想、邓小平理论、“三个代表”重要思想、科学发展观的继承和发展，是当代中国马克思主义、二十一世纪马克思主义，是中华文化和中国精神的时代精华，是党和人民实践经验和集体智慧的结晶，是中国特色社会主义理论体系的重要组成部分，是全党全国人民为实现中华民族伟大复兴而奋斗的行动指南，必须长期坚持并不断发展。在习近平新时代中国特色社会主义思想指导下，中国共产党领导全国各族人民，统揽伟大斗争、伟大工程、伟大事业、伟大梦想，推动中国特色社会主义进入了新时代，实现第一个百年奋斗目标，开启了实现第二个百年奋斗目标新征程。
\item 改革开放以来我们取得一切成绩和进步的根本原因，归结起来就是：开辟了中国特色社会主义道路，形成了中国特色社会主义理论体系，确立了中国特色社会主义制度，发展了中国特色社会主义文化。全党同志要倍加珍惜、长期坚持和不断发展党历经艰辛开创的这条道路、这个理论体系、这个制度、这个文化，高举中国特色社会主义伟大旗帜，坚定道路自信、理论自信、制度自信、文化自信，发扬斗争精神，增强斗争本领，贯彻党的基本理论、基本路线、基本方略，为实现推进现代化建设、完成祖国统一、维护世界和平与促进共同发展这三大历史任务，实现第二个百年奋斗目标、实现中华民族伟大复兴的中国梦而奋斗。
\item 中国共产党自成立以来，始终把为中国人民谋幸福、为中华民族谋复兴作为自己的初心使命，历经百年奋斗，从根本上改变了中国人民的前途命运，开辟了实现中华民族伟大复兴的正确道路，展示了马克思主义的强大生命力，深刻影响了世界历史进程，锻造了走在时代前列的中国共产党。经过长期实践，积累了坚持党的领导、坚持人民至上、坚持理论创新、坚持独立自主、坚持中国道路、坚持胸怀天下、坚持开拓创新、坚持敢于斗争、坚持统一战线、坚持自我革命的宝贵历史经验，这是党和人民共同创造的精神财富，必须倍加珍惜、长期坚持，并在实践中不断丰富和发展。
\item 我国的社会主义建设，必须从我国的国情出发，走中国特色社会主义道路，以中国式现代化全面推进中华民族伟大复兴。在现阶段，我国社会的主要矛盾是人民日益增长的美好生活需要和不平衡不充分的发展之间的矛盾。
\item 我国社会主义建设的根本任务，是进一步解放生产力，发展生产力，逐步实现社会主义现代化，并且为此而改革生产关系和上层建筑中不适应生产力发展的方面和环节。
\item 必须坚持和完善公有制为主体、多种所有制经济共同发展，按劳分配为主体、多种分配方式并存，社会主义市场经济体制等基本经济制度，鼓励一部分地区和一部分人先富起来，逐步实现全体人民共同富裕，在生产发展和社会财富增长的基础上不断满足人民日益增长的美好生活需要，促进人的全面发展。
\item 发展是中国共产党执政兴国的第一要务。必须坚持以人民为中心的发展思想，把握新发展阶段，贯彻创新、协调、绿色、开放、共享的新发展理念，加快构建以国内大循环为主体、国内国际双循环相互促进的新发展格局，推动高质量发展。各项工作都要把有利于发展社会主义社会的生产力，有利于增强社会主义国家的综合国力，有利于提高人民的生活水平，作为总的出发点和检验标准，尊重劳动、尊重知识、尊重人才、尊重创造，做到发展为了人民、发展依靠人民、发展成果由人民共享。
\item 中国共产党在领导社会主义事业中，必须坚持以经济建设为中心，其他各项工作都服从和服务于这个中心。要实施科教兴国战略、人才强国战略、创新驱动发展战略、乡村振兴战略、区域协调发展战略、可持续发展战略、军民融合发展战略，充分发挥科学技术作为第一生产力的作用，充分发挥人才作为第一资源的作用，充分发挥创新作为引领发展第一动力的作用，依靠科技进步，提高劳动者素质，促进国民经济更高质量、更有效率、更加公平、更可持续、更为安全发展。
\item 坚持党的领导、人民当家作主、依法治国有机统一，走中国特色社会主义政治发展道路、中国特色社会主义法治道路，扩大社会主义民主，建设中国特色社会主义法治体系，建设社会主义法治国家，巩固人民民主专政，建设社会主义政治文明。坚持和完善人民代表大会制度、中国共产党领导的多党合作和政治协商制度、民族区域自治制度以及基层群众自治制度。发展更加广泛、更加充分、更加健全的全过程人民民主，推进协商民主广泛多层制度化发展，切实保障人民管理国家事务和社会事务、管理经济和文化事业的权利。尊重和保障人权。广开言路，建立健全民主选举、民主协商、民主决策、民主管理、民主监督的制度和程序。完善中国特色社会主义法律体系，加强法律实施工作，实现国家各项工作法治化。
\item 坚持总体国家安全观，统筹发展和安全，坚决维护国家主权、安全、发展利益。
\item 中国共产党坚持对人民解放军和其他人民武装力量的绝对领导，贯彻习近平强军思想，加强人民解放军的建设，坚持政治建军、改革强军、科技强军、人才强军、依法治军，建设一支听党指挥、能打胜仗、作风优良的人民军队，把人民军队建设成为世界一流军队，切实保证人民解放军有效履行新时代军队使命任务，充分发挥人民解放军在巩固国防、保卫祖国和参加社会主义现代化建设中的作用。
\item 全面准确、坚定不移贯彻“一个国家、两种制度”的方针，促进香港、澳门长期繁荣稳定，坚决反对和遏制“台独”，完成祖国统一大业。
\item 在国际事务中，弘扬和平、发展、公平、正义、民主、自由的全人类共同价值，坚持正确义利观，维护我国的独立和主权，反对霸权主义和强权政治，维护世界和平，促进人类进步，推动构建人类命运共同体，推动建设持久和平、普遍安全、共同繁荣、开放包容、清洁美丽的世界。
\item 中国共产党要领导全国各族人民实现第二个百年奋斗目标、实现中华民族伟大复兴的中国梦，必须紧密围绕党的基本路线，坚持和加强党的全面领导，坚持党要管党、全面从严治党，弘扬坚持真理、坚守理想，践行初心、担当使命，不怕牺牲、英勇斗争，对党忠诚、不负人民的伟大建党精神，加强党的长期执政能力建设、先进性和纯洁性建设，以改革创新精神全面推进党的建设新的伟大工程，以党的政治建设为统领，全面推进党的政治建设、思想建设、组织建设、作风建设、纪律建设，把制度建设贯穿其中，深入推进反腐败斗争，全面提高党的建设科学化水平，以伟大自我革命引领伟大社会革命。
\item 坚持立党为公、执政为民，发扬党的优良传统和作风，不断提高党的领导水平和执政水平，提高拒腐防变和抵御风险的能力，不断增强自我净化、自我完善、自我革新、自我提高能力，不断增强党的阶级基础和扩大党的群众基础，不断提高党的创造力、凝聚力、战斗力，建设学习型、服务型、创新型的马克思主义执政党，使中国共产党始终走在时代前列，成为领导全国人民沿着中国特色社会主义道路不断前进的坚强核心。
\item 党的建设必须坚决实现以下六项基本要求：第一，坚持党的基本路线；第二，坚持解放思想，实事求是，与时俱进，求真务实；第三，坚持新时代党的组织路线；第四，坚持全心全意为人民服务；第五，坚持民主集中制；第六，坚持从严管党治党。
\item 必须提高政治判断力、政治领悟力、政治执行力，增强贯彻落实党的理论和路线方针政策的自觉性和坚定性。
\item 党的思想路线是一切从实际出发，理论联系实际，实事求是，在实践中检验真理和发展真理。全党必须坚持这条思想路线，积极探索，大胆试验，开拓创新，创造性地开展工作，不断研究新情况，总结新经验，解决新问题，在实践中丰富和发展马克思主义，推进马克思主义中国化时代化。
\item 全面贯彻习近平新时代中国特色社会主义思想，以组织体系建设为重点，着力培养忠诚干净担当的高素质干部，着力集聚爱国奉献的各方面优秀人才，坚持德才兼备、以德为先、任人唯贤，为坚持和加强党的全面领导、坚持和发展中国特色社会主义提供坚强组织保证。全党必须增强党组织的政治功能和组织功能，培养选拔党和人民需要的好干部，培养和造就大批堪当时代重任的社会主义事业接班人，聚天下英才而用之，从组织上保证党的基本理论、基本路线、基本方略的贯彻落实。
\item 中国共产党的最大政治优势是密切联系群众，党执政后的最大危险是脱离群众。党风问题、党同人民群众联系问题是关系党生死存亡的问题。
\item 全面从严治党永远在路上，党的自我革命永远在路上。
\item 中国共产党的领导是中国特色社会主义最本质的特征，是中国特色社会主义制度的最大优势，党是最高政治领导力量。
\item 年满十八岁的中国工人、农民、军人、知识分子和其他社会阶层的先进分子，承认党的纲领和章程，愿意参加党的一个组织并在其中积极工作、执行党的决议和按期交纳党费的，可以申请加入中国共产党。
\item 中国共产党党员是中国工人阶级的有共产主义觉悟的先锋战士。中国共产党党员必须全心全意为人民服务，不惜牺牲个人的一切，为实现共产主义奋斗终身。中国共产党党员永远是劳动人民的普通一员。除了法律和政策规定范围内的个人利益和工作职权以外，所有共产党员都不得谋求任何私利和特权。
\item 预备党员的预备期为一年。党组织对预备党员应当认真教育和考察。预备党员的义务同正式党员一样。预备党员的权利，除了没有表决权、选举权和被选举权以外，也同正式党员一样。
\item 预备党员的预备期，从支部大会通过他为预备党员之日算起。党员的党龄，从预备期满转为正式党员之日算起。
\item 党的基层组织是党在社会基层组织中的战斗堡垒，是党的全部工作和战斗力的基础。
\item 对党员的纪律处分有五种：警告、严重警告、撤销党内职务、留党察看、开除党籍。
\item 中国共产党党徽为镰刀和锤头组成的图案。中国共产党党旗为旗面缀有金黄色党徽图案的红旗。中国共产党的党徽党旗是中国共产党的象征和标志。党的各级组织和每一个党员都要维护党徽党旗的尊严。要按照规定制作和使用党徽党旗。
\end{enumerate}

\noindent\textbf{学习参考建议：}
\begin{itemize}[leftmargin=2em,label={}]
\item 《中国共产党章程》（中国共产党第二十次全国代表大会部分修改，2022年10月22日通过）。
\item “学习强国”平台、共产党员网平台上关于马克思主义理论、中国共产党党史、时政要闻等方面的内容。
\end{itemize}

\subsection*{三、重要党史时间节点（示例）}

\begin{enumerate}[label=\arabic*., leftmargin=2.3em]
\item 1921年7月，中国共产党正式宣告成立。
\item 1935年1月，中共中央政治局在长征途中举行的遵义会议，实际上确立了以毛泽东同志为代表的新的中央的正确领导，在党的历史上是一个生死攸关的转折点。
\item 1949年10月1日，中华人民共和国成立。到1956年，中国共产党领导全国各族人民迅速恢复了国民经济并开展了有计划的经济建设，在全国绝大部分地区基本上完成了“三大改造”，确立了生产资料公有制，实现了新民主主义到社会主义的转变。
\item 1978年12月18日至22日，中共十一届三中全会在北京召开。全会冲破长期“左”的错误思想严重束缚，彻底否定了“两个凡是”的错误方针，果断停止使用“以阶级斗争为纲”的口号，作出了把工作重点转移到社会主义现代化建设上来和实行改革开放的战略决策。会后，形成了以邓小平为核心的党中央领导集体，揭开了改革开放的序幕。以这次全会为标志，中国进入了改革开放和社会主义现代化建设的历史新时期。
\item 2022年10月16日，中国共产党第二十次全国代表大会开幕。习近平代表第十九届中央委员会向大会作了题为《高举中国特色社会主义伟大旗帜 为全面建设社会主义现代化国家而团结奋斗》的报告。中国共产党第二十次全国代表大会，是在全党全国各族人民迈上全面建设社会主义现代化国家新征程、向第二个百年奋斗目标进军的关键时刻召开的一次十分重要的大会。大会的主题是：高举中国特色社会主义伟大旗帜，全面贯彻新时代中国特色社会主义思想，弘扬伟大建党精神，自信自强、守正创新，踔厉奋发、勇毅前行，为全面建设社会主义现代化国家、全面推进中华民族伟大复兴而团结奋斗。
\end{enumerate}

\noindent\textbf{学习参考建议：}
\begin{itemize}[leftmargin=2em,label={}]
\item 《中国共产党简史》，本书编写组，人民出版社，2021年2月第1版。
\end{itemize}

\section*{第二部分}

\subsection*{一、《国际歌》第一段歌词}

起来，饥寒交迫的奴隶！起来，全世界受苦的人！满腔的热血已经沸腾，要为真理而斗争！旧世界打个落花流水，奴隶们起来，起来！不要说我们一无所有，我们要做天下的主人！这是最后的斗争，团结起来到明天，英特纳雄耐尔就一定要实现！这是最后的斗争，团结起来到明天，英特纳雄耐尔就一定要实现！

\subsection*{二、中国共产党在社会主义初级阶段的基本路线}

领导和团结全国各族人民，以经济建设为中心，坚持四项基本原则，坚持改革开放，自力更生，艰苦创业，为把我国建设成为富强民主文明和谐美丽的社会主义现代化强国而奋斗。

\subsection*{三、党的性质}

中国共产党是中国工人阶级的先锋队，同时是中国人民和中华民族的先锋队，是中国特色社会主义事业的领导核心，代表中国先进生产力的发展要求，代表中国先进文化的前进方向，代表中国最广大人民的根本利益。党的最高理想和最终目标是实现共产主义。

\subsection*{四、入党誓词}

我志愿加入中国共产党，拥护党的纲领，遵守党的章程，履行党员义务，执行党的决定，严守党的纪律，保守党的秘密，对党忠诚，积极工作，为共产主义奋斗终身，随时准备为党和人民牺牲一切，永不叛党。

\subsection*{五、党员义务}

\noindent （一）认真学习马克思列宁主义、毛泽东思想、邓小平理论、“三个代表”重要思想、科学发展观、习近平新时代中国特色社会主义思想，学习党的路线、方针、政策和决议，学习党的基本知识和党的历史，学习科学、文化、法律和业务知识，努力提高为人民服务的本领。

\noindent （二）增强“四个意识”、坚定“四个自信”、做到“两个维护”，贯彻执行党的基本路线和各项方针、政策，带头参加改革开放和社会主义现代化建设，带动群众为经济发展和社会进步艰苦奋斗，在生产、工作、学习和社会生活中起先锋模范作用。

\noindent （三）坚持党和人民的利益高于一切，个人利益服从党和人民的利益，吃苦在前，享受在后，克己奉公，多做贡献。

\noindent （四）自觉遵守党的纪律，首先是党的政治纪律和政治规矩，模范遵守国家的法律法规，严格保守党和国家的秘密，执行党的决定，服从组织分配，积极完成党的任务。

\noindent （五）维护党的团结和统一，对党忠诚老实，言行一致，坚决反对一切派别组织和小集团活动，反对阳奉阴违的两面派行为和一切阴谋诡计。

\noindent （六）切实开展批评和自我批评，勇于揭露和纠正违反党的原则的言行和工作中的缺点、错误，坚决同消极腐败现象作斗争。

\noindent （七）密切联系群众，向群众宣传党的主张，遇事同群众商量，及时向党反映群众的意见和要求，维护群众的正当利益。

\noindent （八）发扬社会主义新风尚，带头实践社会主义核心价值观和社会主义荣辱观，提倡共产主义道德，弘扬中华民族传统美德，为了保护国家和人民的利益，在一切困难和危险的时刻挺身而出，英勇斗争，不怕牺牲。

\subsection*{六、党员权利}

\noindent （一）参加党的有关会议，阅读党的有关文件，接受党的教育和培训。

\noindent （二）在党的会议上和党报党刊上，参加关于党的政策问题的讨论。

\noindent （三）对党的工作提出建议和倡议。

\noindent （四）在党的会议上有根据地批评党的任何组织和任何党员，向党负责地揭发、检举党的任何组织和任何党员违法乱纪的事实，要求处分违法乱纪的党员，要求罢免或撤换不称职的干部。

\noindent （五）行使表决权、选举权，有被选举权。

\noindent （六）在党组织讨论决定对党员的党纪处分或作出鉴定时，本人有权参加和进行申辩，其他党员可以为他作证和辩护。

\noindent （七）对党的决议和政策如有不同意见，在坚决执行的前提下，可以声明保留，并且可以把自己的意见向党的上级组织直至中央提出。

\noindent （八）向党的上级组织直至中央提出请求、申诉和控告，并要求有关组织给以负责的答复。

党的任何一级组织直至中央都无权剥夺党员的上述权利。

\subsection*{七、党的民主集中制的基本原则}

\noindent （一）党员个人服从党的组织，少数服从多数，下级组织服从上级组织，全党各个组织和全体党员服从党的全国代表大会和中央委员会。

\noindent （二）党的各级领导机关，除它们派出的代表机关和在非党组织中的党组外，都由选举产生。

\noindent （三）党的最高领导机关，是党的全国代表大会和它所产生的中央委员会。党的地方各级领导机关，是党的地方各级代表大会和它们所产生的委员会。党的各级委员会向同级的代表大会负责并报告工作。

\noindent （四）党的上级组织要经常听取下级组织和党员群众的意见，及时解决他们提出的问题。党的下级组织既要向上级组织请示和报告工作，又要独立负责地解决自己职责范围内的问题。上下级组织之间要互通情报、互相支持和互相监督。党的各级组织要按规定实行党务公开，使党员对党内事务有更多的了解和参与。

\noindent （五）党的各级委员会实行集体领导和个人分工负责相结合的制度。凡属重大问题都要按照集体领导、民主集中、个别酝酿、会议决定的原则，由党的委员会集体讨论，作出决定；委员会成员要根据集体的决定和分工，切实履行自己的职责。

\noindent （六）党禁止任何形式的个人崇拜。要保证党的领导人的活动处于党和人民的监督之下，同时维护一切代表党和人民利益的领导人的威信。

\subsection*{八、改革开放以来我们取得一切成绩和进步的根本原因}

开辟了中国特色社会主义道路，形成了中国特色社会主义理论体系，确立了中国特色社会主义制度，发展了中国特色社会主义文化。

\subsection*{九、中国特色社会主义道路、制度和理论体系的关系}

中国特色社会主义道路是实现途径，中国特色社会主义理论体系是行动指南，中国特色社会主义制度是根本保障，三者统一于中国特色社会主义伟大实践，这是党领导人民在建设社会主义长期实践中形成的最鲜明特色。

\subsection*{十、毛泽东思想}

毛泽东思想的活的灵魂三个基本方面：实事求是，群众路线，独立自主。三大法宝：武装斗争、统一战线、党的建设。三大优良作风：理论联系实际的作风，密切联系群众的作风，批评与自我批评的作风。

\subsection*{十一、四个意识}

党的十八届六中全会通过的《关于新形势下党内政治生活的若干准则》强调，全党必须牢固树立政治意识、大局意识、核心意识、看齐意识，自觉在思想上政治上行动上同党中央保持高度一致。

\subsection*{十二、四个自信}

“四个自信”即中国特色社会主义道路自信、理论自信、制度自信和文化自信，由习近平总书记在庆祝中国共产党成立95周年大会上提出，是对党的十八大提出的中国特色社会主义“三个自信”的创造性拓展和完善。

\subsection*{十三、两个维护}

坚决维护习近平总书记党中央的核心、全党的核心地位，坚决维护党中央权威和集中统一领导。坚决做到“两个维护”，既是根本政治任务，也是根本政治纪律和政治规矩，是牢固树立“四个意识”的集中体现。

\subsection*{十四、“五位一体”总体布局}

“五位一体”总体布局是指经济建设、政治建设、文化建设、社会建设和生态文明建设五位一体，全面推进。

\subsection*{十五、“四个全面”战略布局}

“四个全面”战略布局是指全面建设社会主义现代化国家、全面深化改革、全面依法治国、全面从严治党。

\subsection*{十六、三个务必}

中国共产党已走过百年奋斗历程。中国共产党立志于中华民族千秋伟业，致力于人类和平与发展崇高事业，责任无比重大，使命无上光荣。全党同志务必不忘初心、牢记使命，务必谦虚谨慎、艰苦奋斗，务必敢于斗争、善于斗争，坚定历史自信，增强历史主动，谱写新时代中国特色社会主义更加绚丽的华章。

\subsection*{十七、十年来三件大事}

一是迎来中国共产党成立一百周年，二是中国特色社会主义进入新时代，三是完成脱贫攻坚、全面建成小康社会的历史任务，实现第一个百年奋斗目标。这是中国共产党和中国人民团结奋斗赢得的历史性胜利，是彪炳中华民族发展史册的历史性胜利，也是对世界具有深远影响的历史性胜利。

\subsection*{十八、跳出历史周期率的第二个答案}

经过不懈努力，党找到了自我革命这一跳出治乱兴衰历史周期率的第二个答案，确保党永远不变质、不变色、不变味。

\subsection*{十九、归根到底是两个“行”}

马克思主义是我们立党立国、兴党兴国的根本指导思想。实践告诉我们，中国共产党为什么能，中国特色社会主义为什么好，归根到底是马克思主义行，是中国化时代化的马克思主义行。拥有马克思主义科学理论指导是中国共产党坚定信仰信念、把握历史主动的根本所在。

\subsection*{二十、党的中心任务}

从现在起，中国共产党的中心任务就是团结带领全国各族人民全面建成社会主义现代化强国、实现第二个百年奋斗目标，以中国式现代化全面推进中华民族伟大复兴。

\subsection*{二十一、中国式现代化的五个特色}

中国式现代化，是中国共产党领导的社会主义现代化，既有各国现代化的共同特征，更有基于自己国情的中国特色。中国式现代化是人口规模巨大的现代化，是全体人民共同富裕的现代化，是物质文明和精神文明相协调的现代化，是人与自然和谐共生的现代化，是走和平发展道路的现代化。

\subsection*{二十二、中国式现代化的本质要求}

中国式现代化的本质要求是：坚持中国共产党领导，坚持中国特色社会主义，实现高质量发展，发展全过程人民民主，丰富人民精神世界，实现全体人民共同富裕，促进人与自然和谐共生，推动构建人类命运共同体，创造人类文明新形态。

\subsection*{二十三、五个重大原则}

前进道路上，必须牢牢把握以下重大原则：坚持和加强党的全面领导，坚持中国特色社会主义道路，坚持以人民为中心的发展思想，坚持深化改革开放，坚持发扬斗争精神。要增强全党全国各族人民的志气、骨气、底气，不信邪、不怕鬼、不怕压，知难而进、迎难而上，统筹发展和安全，全力战胜前进道路上各种困难和挑战，依靠顽强斗争打开事业发展新天地。

\subsection*{二十四、全面建设社会主义现代化国家的首要任务}

高质量发展是全面建设社会主义现代化国家的首要任务。发展是党执政兴国的第一要务。没有坚实的物质技术基础，就不可能全面建成社会主义现代化强国。必须完整、准确、全面贯彻新发展理念，坚持社会主义市场经济改革方向，坚持高水平对外开放，加快构建以国内大循环为主体、国内国际双循环相互促进的新发展格局。

\subsection*{二十五、全面建设社会主义现代化国家的基础性、战略性支撑}

教育、科技、人才是全面建设社会主义现代化国家的基础性、战略性支撑。必须坚持科技是第一生产力、人才是第一资源、创新是第一动力，深入实施科教兴国战略、人才强国战略、创新驱动发展战略，开辟发展新领域新赛道，不断塑造发展新动能新优势。

\subsection*{二十六、全面建设社会主义现代化国家的应有之义}

我国是工人阶级领导的、以工农联盟为基础的人民民主专政的社会主义国家，国家一切权力属于人民。人民民主是社会主义的生命，是全面建设社会主义现代化国家的应有之义。全过程人民民主是社会主义民主政治的本质属性，是最广泛、最真实、最管用的民主。必须坚定不移走中国特色社会主义政治发展道路，坚持党的领导、人民当家作主、依法治国有机统一，坚持人民主体地位，充分体现人民意志、保障人民权益、激发人民创造活力。

\subsection*{二十七、在法治轨道上全面建设社会主义现代化国家}

全面依法治国是国家治理的一场深刻革命，关系党执政兴国，关系人民幸福安康，关系党和国家长治久安。必须更好发挥法治固根本、稳预期、利长远的保障作用，在法治轨道上全面建设社会主义现代化国家。

\subsection*{二十八、全面建设社会主义现代化国家必须坚持中国特色社会主义文化发展道路}

全面建设社会主义现代化国家，必须坚持中国特色社会主义文化发展道路，增强文化自信，围绕举旗帜、聚民心、育新人、兴文化、展形象建设社会主义文化强国，发展面向现代化、面向世界、面向未来的，民族的科学的大众的社会主义文化，激发全民族文化创新创造活力，增强实现中华民族伟大复兴的精神力量。

\subsection*{二十九、为民造福是立党为公、执政为民的本质要求}

江山就是人民，人民就是江山。中国共产党领导人民打江山、守江山，守的是人民的心。治国有常，利民为本。为民造福是立党为公、执政为民的本质要求。必须坚持在发展中保障和改善民生，鼓励共同奋斗创造美好生活，不断实现人民对美好生活的向往。

\subsection*{三十、理想信念重要性}

事实一再表明，理想信念动摇是最危险的动摇，理想信念滑坡是最危险的滑坡。……在革命、建设、改革各个历史时期，有无数共产党员为了党和人民事业英勇牺牲了，支撑他们的就是“革命理想高于天”的精神力量。

\subsection*{三十一、全面建设社会主义现代化国家的内在要求}

尊重自然、顺应自然、保护自然，是全面建设社会主义现代化国家的内在要求。必须牢固树立和践行绿水青山就是金山银山的理念，站在人与自然和谐共生的高度谋划发展。

\subsection*{三十二、国家安全是民族复兴的根基}

国家安全是民族复兴的根基，社会稳定是国家强盛的前提。必须坚定不移贯彻总体国家安全观，把维护国家安全贯穿党和国家工作各方面全过程，确保国家安全和社会稳定。

\subsection*{三十三、全面建设社会主义现代化国家的战略要求}

如期实现建军一百年奋斗目标，加快把人民军队建成世界一流军队，是全面建设社会主义现代化国家的战略要求。必须贯彻新时代党的强军思想，贯彻新时代军事战略方针，坚持党对人民军队的绝对领导，坚持政治建军、改革强军、科技强军、人才强军、依法治军，加快军事理论现代化、军队组织形态现代化、军事人员现代化、武器装备现代化，提高捍卫国家主权、安全、发展利益战略能力，有效履行新时代人民军队使命任务。

\subsection*{三十四、最佳制度安排}

“一国两制”是中国特色社会主义的伟大创举，是香港、澳门回归后保持长期繁荣稳定的最佳制度安排，必须长期坚持。

\subsection*{三十五、构建人类命运共同体}

2013年3月，中国国家主席习近平提出推动构建人类命运共同体。人类命运共同体，就是每个国家的前途命运都紧密相连，不同群体风雨同舟、和谐共生、合作共赢。党的二十大报告鲜明指出：“构建人类命运共同体是世界各国人民前途所在。万物并育而不相害，道并行而不相悖。只有各国行天下之大道，和睦相处、合作共赢，繁荣才能持久，安全才有保障。”人类命运共同体理念，蕴含有开放包容、公平正义、和谐共处、多元互鉴、团结协作的丰富内涵。构建人类命运共同体是习近平总书记提出的重大原创性思想，是习近平外交思想的核心理念，是对“建设一个什么样的世界、如何建设这个世界”这一时代课题给出的中国答案，具有宝贵理论价值、重大时代意义和深远世界影响。

\subsection*{三十六、全面建设社会主义现代化国家、全面推进中华民族伟大复兴关键在党}

全面建设社会主义现代化国家、全面推进中华民族伟大复兴，关键在党。这一重要论断，明确了党在新时代中国特色社会主义事业中的领导核心地位。第一，中国共产党的领导核心地位是党章和宪法予以明确的。第二，党史和新中国史充分证明，只有中国共产党才能担当实现中华民族伟大复兴的重任。第三，应对前进道路上的各种风险挑战，必须坚持和加强党的全面领导，充分发挥党员的先锋模范作用。中国共产党作为世界上最大的马克思主义执政党，要始终赢得人民拥护、巩固长期执政地位，必须时刻保持解决大党独有难题的清醒和坚定。全党必须牢记，全面从严治党永远在路上，党的自我革命永远在路上，决不能有松劲歇脚、疲劳厌战的情绪，必须持之以恒推进全面从严治党，深入推进新时代党的建设新的伟大工程，以党的自我革命引领社会革命。

\subsection*{三十七、五个“必由之路”}

全党必须牢记，坚持党的全面领导是坚持和发展中国特色社会主义的必由之路，中国特色社会主义是实现中华民族伟大复兴的必由之路，团结奋斗是中国人民创造历史伟业的必由之路，贯彻新发展理念是新时代我国发展壮大的必由之路，全面从严治党是党永葆生机活力、走好新的赶考之路的必由之路。这是我们在长期实践中得出的至关紧要的规律性认识，必须倍加珍惜、始终坚持，咬定青山不放松，引领和保障中国特色社会主义巍巍巨轮乘风破浪、行稳致远。

\section*{第三部分}

\begin{enumerate}[label=\textbf{材料示例\chinese*}, leftmargin=5.2em]
\item \href{https://www.12371.cn/2025/10/28/ARTI1761640401107119.shtml}{中共中央关于制定国民经济和社会发展第十五个五年规划的建议} \textcolor{gray}{（共产党员网）}
\item \href{https://www.qstheory.cn/20251129/5c573dc1d6574c5b8dd814e5403ae88e/c.html}{推进党的自我革命要做到“五个进一步到位”} \textcolor{gray}{（求是网）}
\item \href{https://mp.weixin.qq.com/s/XLDi-vArQj-8JXKCRT-6ZQ}{中国始终是世界经济发展的积极贡献者和机遇提供者} \textcolor{gray}{（原文链接）}
\item \href{https://mp.weixin.qq.com/s/1b_QMoc5TUANlYwCo6RIPA}{文平：认清“伪史论”的历史虚无主义本质} \textcolor{gray}{（原文链接）}
\item \href{https://www.12371.cn/2026/02/27/ARTI1772153945397806.shtml}{为人民出政绩 以实干出政绩——习近平总书记引领全党以正确政绩观干事创业} \textcolor{gray}{（共产党员网）}
\item \href{https://www.qstheory.cn/20260227/c77ea2492a464a9081e2330e2e906b6f/c.html}{深刻领悟习近平外交思想源于时代引领时代的理论品格} \textcolor{gray}{（求是网）}
\item \href{https://www.12371.cn/2026/01/31/ARTI1769849380447546.shtml}{习近平：发挥比较优势 坚持稳中求进 推动我国未来产业发展不断取得新突破} \textcolor{gray}{（共产党员网）}
\item \href{https://www.12371.cn/2026/02/15/ARTI1771139538258762.shtml}{习近平：当前经济工作的重点任务} \textcolor{gray}{（共产党员网）}
\item \href{https://www.12371.cn/2026/01/12/ARTI1768214272707593.shtml}{习近平在二十届中央纪委五次全会上发表重要讲话强调 以更高标准更实举措推进全面从严治党 为实现“十五五”时期目标任务提供坚强保障} \textcolor{gray}{（共产党员网）}
\end{enumerate}

\section*{练习题}

\subsection*{一、选择题}

\begin{enumerate}[label=\textbf{\arabic*.}, leftmargin=2.5em, itemsep=0.7em]
\item 有关全国性的重大政策问题，只有党中央有权作出决定，各部门、各地方的党组织可以向中央提出\allowbreak\underline{\hspace{4em}}，但不得擅自作出决定和对外发表主张。
\begin{enumerate}[label=\Alph*., leftmargin=2em, itemsep=0.1em, topsep=0.15em]
\item 建议
\item 申诉
\item 要求
\item 抗议
\end{enumerate}
\textbf{【答案】A}
\item 党的下级组织必须坚决执行上级组织的决定。下级组织如果认为上级组织的决定不符合本地区、本部门的实际情况，可以\allowbreak\underline{\hspace{4em}}。
\begin{enumerate}[label=\Alph*., leftmargin=2em, itemsep=0.1em, topsep=0.15em]
\item 请求改变
\item 不予执行
\item 暂缓执行
\item 变通执行
\end{enumerate}
\textbf{【答案】A}
\item 下列说法中，符合《中国共产党章程》规定的是\allowbreak\underline{\hspace{4em}}。
\begin{enumerate}[label=\Alph*., leftmargin=2em, itemsep=0.1em, topsep=0.15em]
\item 党的各级组织的报刊和其他宣传工具，可以有选择地宣传党的纲领、路线、方针和政策。
\item 党的各级组织的报刊和其他宣传工具，必须宣传党的路线、方针、政策和决议。
\item 党的各级组织的主要报刊和其他宣传工具，不必宣传党的纲领、路线、方针和政策。
\item 党的各级组织的报刊和其他宣传工具，自行宣传党的纲领、路线、方针和政策。
\end{enumerate}
\textbf{【答案】B}
\item 党组织讨论决定问题，必须执行\allowbreak\underline{\hspace{4em}}的原则。
\begin{enumerate}[label=\Alph*., leftmargin=2em, itemsep=0.1em, topsep=0.15em]
\item 个人服从组织
\item 少数服从多数
\item 下级服从上级
\item 集体利益优先
\end{enumerate}
\textbf{【答案】B}
\item 党的中央和地方各级委员会可以派出\allowbreak\underline{\hspace{4em}}。
\begin{enumerate}[label=\Alph*., leftmargin=2em, itemsep=0.1em, topsep=0.15em]
\item 办事机构
\item 代表机关
\item 全权大使
\item 领导小组
\end{enumerate}
\textbf{【答案】B}
\item 以下关于巡视制度的说法中，正确的是\allowbreak\underline{\hspace{4em}}。
\begin{enumerate}[label=\Alph*., leftmargin=2em, itemsep=0.1em, topsep=0.15em]
\item 党的中央和省、自治区、直辖市委员会实行巡视制度，在一届任期内，对所管理的地方、部门、企事业单位党组织实现派驻全覆盖。
\item 中央有关部委和国家机关部门党组（党委）根据工作需要，开展巡视工作。
\item 党的市（地、州、盟）和县（市、区、旗）委员会建立巡视制度。
\item 巡视制度是中国共产党治国的一项重要制度。
\end{enumerate}
\textbf{【答案】B}
\item 凡属应由下级组织处理的问题，如无特殊情况，上级领导机关\allowbreak\underline{\hspace{4em}}。
\begin{enumerate}[label=\Alph*., leftmargin=2em, itemsep=0.1em, topsep=0.15em]
\item 可以适当干预
\item 可以代为处理
\item 不要干预
\item 应当干预
\end{enumerate}
\textbf{【答案】C}
\item 《中国共产党章程》规定，党是根据自己的\allowbreak\underline{\hspace{4em}}，按照民主集中制组织起来的统一整体。
\begin{enumerate}[label=\Alph*., leftmargin=2em, itemsep=0.1em, topsep=0.15em]
\item 理想和信念
\item 目标和任务
\item 纲领和章程
\item 行动纲领
\end{enumerate}
\textbf{【答案】C}
\item 《中国共产党章程》规定，党支部是党的\allowbreak\underline{\hspace{4em}}，担负直接教育党员、管理党员、监督党员和组织群众、宣传群众、凝聚群众、服务群众的职责。
\begin{enumerate}[label=\Alph*., leftmargin=2em, itemsep=0.1em, topsep=0.15em]
\item 领导核心
\item 政治核心
\item 坚强核心
\item 基础组织
\end{enumerate}
\textbf{【答案】D}
\item 《中国共产党章程》规定，党的中央和地方纪律检查委员会向同级党和国家机关全面派驻党的纪律检查组。纪律检查组组长\allowbreak\underline{\hspace{4em}}驻在单位党的领导组织的有关会议。他们的工作必须受到该单位党的领导组织的支持。
\begin{enumerate}[label=\Alph*., leftmargin=2em, itemsep=0.1em, topsep=0.15em]
\item 出席
\item 列席
\item 参加
\item 旁听
\end{enumerate}
\textbf{【答案】C}
\item 《中国共产党章程》规定，中国共产党党员永远是劳动人民的普通一员。除了法律和政策规定范围内的个人利益和工作职权以外，所有共产党员都不得谋求任何\allowbreak\underline{\hspace{4em}}。
\begin{enumerate}[label=\Alph*., leftmargin=2em, itemsep=0.1em, topsep=0.15em]
\item 私利和特权
\item 政治权力
\item 经济利益
\item 物质享受
\end{enumerate}
\textbf{【答案】A}
\item 《中国共产党章程》规定，党的民主集中制的基本原则之一是“四个服从”，即党员个人服从党的组织，\allowbreak\underline{\hspace{4em}}，下级组织服从上级组织，全党各个组织和全体党员服从党的全国代表大会和中央委员会。
\begin{enumerate}[label=\Alph*., leftmargin=2em, itemsep=0.1em, topsep=0.15em]
\item 委员服从党委书记
\item 少数服从多数
\item 普通党员服从党组织书记
\item 以上都不是
\end{enumerate}
\textbf{【答案】B}
\item 《中国共产党章程》规定，凡属重大问题，党的各级委员会都要按照集体领导、民主集中、个别酝酿、\allowbreak\underline{\hspace{4em}}的原则，由党的委员会集体讨论，作出决定。
\begin{enumerate}[label=\Alph*., leftmargin=2em, itemsep=0.1em, topsep=0.15em]
\item 会议决定
\item 书记个人决定
\item 常委会决定
\item 书记办公会
\end{enumerate}
\textbf{【答案】A}
\item 党员必须坚持\allowbreak\underline{\hspace{4em}}高于一切，个人利益服从党和人民的利益，吃苦在前，享受在后，克己奉公，多做贡献。
\begin{enumerate}[label=\Alph*., leftmargin=2em, itemsep=0.1em, topsep=0.15em]
\item 个人利益
\item 局部利益
\item 单位利益
\item 党和人民的利益
\end{enumerate}
\textbf{【答案】D}
\item \allowbreak\underline{\hspace{4em}}，是马克思主义政党同其他政党的根本区别。\\ ①维护社会公平正义；②始终同人民在一起；\\ ③为人民利益而奋斗；④促进社会和谐发展
\begin{enumerate}[label=\Alph*., leftmargin=2em, itemsep=0.1em, topsep=0.15em]
\item ①②
\item ②③
\item ①④
\item ③④
\end{enumerate}
\textbf{【答案】B}
\item 党员除了享有表决权、选举权和被选举权以外，还有权要求\allowbreak\underline{\hspace{4em}}或撤换不称职的干部。
\begin{enumerate}[label=\Alph*., leftmargin=2em, itemsep=0.1em, topsep=0.15em]
\item 开除
\item 罢免
\item 辞退
\item 处分
\end{enumerate}
\textbf{【答案】B}
\item 《中国共产党章程》规定，党员对党组织作出的处分决定不服，可以提出\allowbreak\underline{\hspace{4em}}，有关党组织必须负责处理或者迅速转递，不得扣压。
\begin{enumerate}[label=\Alph*., leftmargin=2em, itemsep=0.1em, topsep=0.15em]
\item 申辩
\item 复议
\item 申诉
\item 诉讼
\end{enumerate}
\textbf{【答案】C}
\item 《中国共产党章程》规定，党的纪律是党的\allowbreak\underline{\hspace{4em}}必须遵守的行为规则，是维护党的团结统一、完成党的任务的保证。
\begin{enumerate}[label=\Alph*., leftmargin=2em, itemsep=0.1em, topsep=0.15em]
\item 基层组织
\item 所有干部
\item 各级组织和全体党员
\item 正式党员
\end{enumerate}
\textbf{【答案】C}
\item 中国共产党坚持对人民解放军和其他人民武装力量的绝对领导，贯彻习近平强军思想，加强人民解放军的建设，坚持政治建军、改革强军、\allowbreak\underline{\hspace{4em}}、人才强军、依法治军，建设一支听党指挥、能打胜仗、作风优良的人民军队，把人民军队建设成为世界一流军队
\begin{enumerate}[label=\Alph*., leftmargin=2em, itemsep=0.1em, topsep=0.15em]
\item 科技建军
\item 科技强军
\item 科技兴军
\item 科技为军
\end{enumerate}
\textbf{【答案】B}
\item 全面建设社会主义现代化国家的基础性、战略性支撑是\allowbreak\underline{\hspace{4em}}。
\begin{enumerate}[label=\Alph*., leftmargin=2em, itemsep=0.1em, topsep=0.15em]
\item 教育、科技、人才
\item 教育、科技、创新
\item 科教、人才、创新
\item 科技、人才、创新
\end{enumerate}
\textbf{【答案】A}
\item 党的建设必须坚持党的基本路线。全党要用邓小平理论、“三个代表”重要思想、科学发展观、习近平新时代中国特色社会主义思想和党的基本路线统一思想，统一行动，并且毫不动摇地长期坚持下去。必须把改革开放同\allowbreak\underline{\hspace{4em}}统一起来，全面落实党的基本路线，反对一切“左”的和右的错误倾向，要警惕右，但主要是防止“左”。
\begin{enumerate}[label=\Alph*., leftmargin=2em, itemsep=0.1em, topsep=0.15em]
\item 两个维护
\item 四个自信
\item 四个伟大
\item 四项基本原则
\end{enumerate}
\textbf{【答案】D}
\item \allowbreak\underline{\hspace{4em}}负责召集中央政治局会议和中央政治局常务委员会会议。
\begin{enumerate}[label=\Alph*., leftmargin=2em, itemsep=0.1em, topsep=0.15em]
\item 中央政治局常委
\item 中央书记处书记
\item 中央委员会总书记
\item 中央政治局委员
\end{enumerate}
\textbf{【答案】C}
\item 党的最高领导机关是\allowbreak\underline{\hspace{4em}}。
\begin{enumerate}[label=\Alph*., leftmargin=2em, itemsep=0.1em, topsep=0.15em]
\item 党的全国代表大会和它所产生的中央委员会
\item 党的全国代表大会
\item 党的中央委员会
\item 党的中央政治局
\end{enumerate}
\textbf{【答案】A}
\item 发展党员，必须把\allowbreak\underline{\hspace{4em}}放在首位。
\begin{enumerate}[label=\Alph*., leftmargin=2em, itemsep=0.1em, topsep=0.15em]
\item 政治标准
\item 思想标准
\item 道德标准
\item 行为标准
\end{enumerate}
\textbf{【答案】A}
\item 企业、农村、机关、学校、医院、科研院所、街道社区、社会组织、人民解放军连队和其他基层单位，凡是有正式党员\allowbreak\underline{\hspace{4em}}人及以上，都应当成立党的基层组织。
\begin{enumerate}[label=\Alph*., leftmargin=2em, itemsep=0.1em, topsep=0.15em]
\item 二
\item 三
\item 四
\item 五
\end{enumerate}
\textbf{【答案】B}
\item 中国共产党维护和发展\allowbreak\underline{\hspace{4em}}的社会主义民族关系。
\begin{enumerate}[label=\Alph*., leftmargin=2em, itemsep=0.1em, topsep=0.15em]
\item 平等团结友爱互助
\item 平等友爱和谐互助
\item 平等团结互助和谐
\item 平等团结和睦互助
\end{enumerate}
\textbf{【答案】C}
\item 中国积极参与全球治理体系改革和建设，践行\allowbreak\underline{\hspace{4em}}的全球治理观。
\begin{enumerate}[label=\Alph*., leftmargin=2em, itemsep=0.1em, topsep=0.15em]
\item 共商共建共享
\item 互帮互助互赢
\item 互商互利互赢
\item 共赢共享共利
\end{enumerate}
\textbf{【答案】A}
\item 中国共产党党徽为镰刀和\allowbreak\underline{\hspace{4em}}组成的图案。
\begin{enumerate}[label=\Alph*., leftmargin=2em, itemsep=0.1em, topsep=0.15em]
\item 斧头
\item 锄头
\item 五星
\item 锤头
\end{enumerate}
\textbf{【答案】D}
\item “三会一课”制度的内容是\allowbreak\underline{\hspace{4em}}。
\begin{enumerate}[label=\Alph*., leftmargin=2em, itemsep=0.1em, topsep=0.15em]
\item 支部党员大会，党支部委员会，党小组会，党课
\item 支部党员大会，党支部委员会，领导班子会，党课
\item 民主生活会，党支部委员会，组织生活会，党课
\item 支部党员大会，党支部委员会，组织生活会，党课
\end{enumerate}
\textbf{【答案】A}
\item 中国共产党同全国各民族工人、农民、知识分子团结在一起，同各民主党派、\allowbreak\underline{\hspace{4em}}、各民族的爱国力量团结在一起，组成最广泛的爱国统一战线。
\begin{enumerate}[label=\Alph*., leftmargin=2em, itemsep=0.1em, topsep=0.15em]
\item 无党派人士
\item 无党派民主人士
\item 政协代表
\item 人大代表
\end{enumerate}
\textbf{【答案】A}
\item 健全总揽全局、协调各方的党的领导制度体系，完善党中央重大决策部署落实机制，确保全党在政治立场、\allowbreak\underline{\hspace{4em}}、政治原则、政治道路上同党中央保持高度一致，确保党的团结统一。
\begin{enumerate}[label=\Alph*., leftmargin=2em, itemsep=0.1em, topsep=0.15em]
\item 政治制度
\item 政治选择
\item 政治素养
\item 政治方向
\end{enumerate}
\textbf{【答案】D}
\item 全面建设社会主义现代化国家的首要任务是\allowbreak\underline{\hspace{4em}}。
\begin{enumerate}[label=\Alph*., leftmargin=2em, itemsep=0.1em, topsep=0.15em]
\item 高质量发展
\item 高速发展
\item 高资源利用率
\item 高效创新
\end{enumerate}
\textbf{【答案】A}
\item 中国奉行\allowbreak\underline{\hspace{4em}}的外交政策。
\begin{enumerate}[label=\Alph*., leftmargin=2em, itemsep=0.1em, topsep=0.15em]
\item 独立自主
\item 不受外部势力干涉
\item 符合联合国要求
\item 符合世界人民要求
\end{enumerate}
\textbf{【答案】A}
\item \allowbreak\underline{\hspace{4em}}，这是马克思主义政治经济学的根本立场。
\begin{enumerate}[label=\Alph*., leftmargin=2em, itemsep=0.1em, topsep=0.15em]
\item 发展依靠人民
\item 发展为了人民
\item 推动经济发展
\item 发展才是硬道理
\end{enumerate}
\textbf{【答案】B}
\item 就业是\allowbreak\underline{\hspace{4em}}，要坚持经济发展就业导向，实施就业优先战略和积极就业政策，千方百计稳定和扩大就业，实现\allowbreak\underline{\hspace{4em}}就业。
\begin{enumerate}[label=\Alph*., leftmargin=2em, itemsep=0.1em, topsep=0.15em]
\item 民生之本 更高质量和更有效率
\item 民生之本 更高质量和更充分
\item 经济血脉 更高质量和更有效率
\item 经济血脉 更高质量和更充分
\end{enumerate}
\textbf{【答案】B}
\item 习近平经济思想坚持马克思主义政治经济学的基本原理和方法，在适应新形势、解决新问题、应对新挑战中形成了一系列具有鲜明时代性和创造性的理论成果，集中体现了中国共产党对经济发展规律特别是\allowbreak\underline{\hspace{4em}}的深刻洞见，为丰富发展马克思主义政治经济学作出了重要\allowbreak\underline{\hspace{4em}}贡献。
\begin{enumerate}[label=\Alph*., leftmargin=2em, itemsep=0.1em, topsep=0.15em]
\item 社会主义经济建设规律 集成性
\item 社会主义市场经济规律 集成性
\item 社会主义经济建设规律 原创性
\item 社会主义市场经济规律 原创性
\end{enumerate}
\textbf{【答案】C}
\item 坚持实事求是，必须\allowbreak\underline{\hspace{4em}}。
\begin{enumerate}[label=\Alph*., leftmargin=2em, itemsep=0.1em, topsep=0.15em]
\item 解放思想、恪守传统、求真务实
\item 解放思想、与时俱进、求真务实
\item 恪守传统、与时俱进、求真务实
\item 解放思想、与时俱进、恪守传统
\end{enumerate}
\textbf{【答案】B}
\item \allowbreak\underline{\hspace{4em}}的统一，内在地贯穿于马克思主义哲学的每一条原理、每一个观点之中。
\begin{enumerate}[label=\Alph*., leftmargin=2em, itemsep=0.1em, topsep=0.15em]
\item 唯物主义和历史观
\item 唯物主义和辩证法
\item 唯物主义自然观和历史观
\item 唯物主义自然观和辩证法
\end{enumerate}
\textbf{【答案】B}
\item 从动态看，生产关系是由\allowbreak\underline{\hspace{4em}}四个环节构成的。
\begin{enumerate}[label=\Alph*., leftmargin=2em, itemsep=0.1em, topsep=0.15em]
\item 生产、再生产、交换和消费
\item 生产、分配、交换和消费
\item 生产、分配、再生产和消费
\item 生产、分配、交换和再生产
\end{enumerate}
\textbf{【答案】B}
\item 习近平总书记指出：“必须看到，全面建设社会主义现代化国家，实现中华民族伟大复兴，最艰巨最繁重的任务依然在\allowbreak\underline{\hspace{4em}}，最广泛最深厚的基础依然在\allowbreak\underline{\hspace{4em}}。”
\begin{enumerate}[label=\Alph*., leftmargin=2em, itemsep=0.1em, topsep=0.15em]
\item 乡镇乡镇
\item 农村农村
\item 城市城市
\item 沿海城市沿海城市
\end{enumerate}
\textbf{【答案】B}
\item 党的群众观点主要包括\allowbreak\underline{\hspace{4em}}。\\ ①一切为了人民群众的观点\\ ②一切依靠人民群众的观点\\ ③向人民群众负责的观点\\ ④向人民群众学习的观点
\begin{enumerate}[label=\Alph*., leftmargin=2em, itemsep=0.1em, topsep=0.15em]
\item ①②③
\item ①③④
\item ②③④
\item ①②③④
\end{enumerate}
\textbf{【答案】D}
\item “从群众中来，到群众中去”的群众路线，与“从实践到认识，从认识到实践”的认识路线在本质上是\allowbreak\underline{\hspace{4em}}。二者统一的基础就是\allowbreak\underline{\hspace{4em}}。
\begin{enumerate}[label=\Alph*., leftmargin=2em, itemsep=0.1em, topsep=0.15em]
\item 一致的 实践
\item 一致的 认识
\item 不同的 实践
\item 不同的 认识
\end{enumerate}
\textbf{【答案】A}
\item 人类历史是在社会基本矛盾运动中进行的，而历史的主体是\allowbreak\underline{\hspace{4em}}。
\begin{enumerate}[label=\Alph*., leftmargin=2em, itemsep=0.1em, topsep=0.15em]
\item 历史人物
\item 人民群众
\item 英雄人物
\item 政治家
\end{enumerate}
\textbf{【答案】B}
\item 实践的观点是马克思主义哲学的\allowbreak\underline{\hspace{4em}}。
\begin{enumerate}[label=\Alph*., leftmargin=2em, itemsep=0.1em, topsep=0.15em]
\item 核心观点
\item 基本观点
\item 重要观点
\item 一般观点
\end{enumerate}
\textbf{【答案】A}
\item \allowbreak\underline{\hspace{4em}}是适度原则的体现，体现了度的两个关节点的对立统一，是“底”与“顶”的有机结合。
\begin{enumerate}[label=\Alph*., leftmargin=2em, itemsep=0.1em, topsep=0.15em]
\item 辩证思维
\item 系统思维
\item 底线思维
\item 创新思维
\end{enumerate}
\textbf{【答案】C}
\item 人的正确思想，只能从\allowbreak\underline{\hspace{4em}}来。
\begin{enumerate}[label=\Alph*., leftmargin=2em, itemsep=0.1em, topsep=0.15em]
\item 理论思考中
\item 自己头脑中
\item 社会实践中
\item 天上掉下
\end{enumerate}
\textbf{【答案】C}
\item 加强社会主义核心价值体系建设，坚持马克思主义指导思想，树立中国特色社会主义共同理想，弘扬以爱国主义为核心的民族精神和以\allowbreak\underline{\hspace{4em}}为核心的时代精神。
\begin{enumerate}[label=\Alph*., leftmargin=2em, itemsep=0.1em, topsep=0.15em]
\item 改革创新
\item 爱岗敬业
\item 诚信友善
\item 守正务实
\end{enumerate}
\textbf{【答案】A}
\item 党的地方各级领导机关，是\allowbreak\underline{\hspace{4em}}。
\begin{enumerate}[label=\Alph*., leftmargin=2em, itemsep=0.1em, topsep=0.15em]
\item 党的地方各级代表大会和它们所产生的委员会
\item 党的地方各级委员会
\item 党的地方各级代表大会
\item 党的地方各级代表大会
\end{enumerate}
\textbf{【答案】A}
\item 预备党员的权利，除了没有\allowbreak\underline{\hspace{4em}}、选举权和被选举权以外，同正式党员一样。
\begin{enumerate}[label=\Alph*., leftmargin=2em, itemsep=0.1em, topsep=0.15em]
\item 表决权
\item 控告权
\item 申诉权
\item 批评权
\end{enumerate}
\textbf{【答案】A}
\item 对党的中央委员会和地方各级委员会的委员、候补委员，给以撤销党内职务、留党察看或开除党籍的处分，在\allowbreak\underline{\hspace{4em}}，可以先由中央政治局和地方各级委员会常务委员会作出处理决定，待召开委员会全体会议时予以追认。
\begin{enumerate}[label=\Alph*., leftmargin=2em, itemsep=0.1em, topsep=0.15em]
\item 全体会议闭会期间
\item 参会人数不足时
\item 特殊情况下
\item 其他重要会议期间
\end{enumerate}
\textbf{【答案】A}
\item 全党同志要坚定道路自信、理论自信、制度自信、文化自信，发扬\allowbreak\underline{\hspace{4em}}，增强斗争本领。
\begin{enumerate}[label=\Alph*., leftmargin=2em, itemsep=0.1em, topsep=0.15em]
\item 斗争精神
\item 革命精神
\item 开拓精神
\item 忘我精神
\end{enumerate}
\textbf{【答案】A}
\item 以下关于巡视制度的说法中，正确的是\allowbreak\underline{\hspace{4em}}。
\begin{enumerate}[label=\Alph*., leftmargin=2em, itemsep=0.1em, topsep=0.15em]
\item 党的市（地、州、盟）和县（市、区、旗）委员会建立巡视制度
\item 中央有关部委和国家机关部门党组（党委）根据工作需要，开展巡视工作
\item 党的中央和省、自治区、直辖市委员会实行巡视制度，在一届任期内，对所管理的地方、部门、企事业单位党组织实现派驻全覆盖
\item 中央有关部委和国家机关部门党组（党委）每年一定会开展巡视工作
\end{enumerate}
\textbf{【答案】B}
\item 预备党员的预备期为\allowbreak\underline{\hspace{4em}}。
\begin{enumerate}[label=\Alph*., leftmargin=2em, itemsep=0.1em, topsep=0.15em]
\item 2年
\item 1年
\item 6个月
\item 3个月
\end{enumerate}
\textbf{【答案】B}
\item 党的二十大报告指出，我们要推进\allowbreak\underline{\hspace{4em}}，坚持山水林田湖草沙一体化保护和系统治理，统筹产业结构调整、污染治理、生态保护、应对气候变化，协同推进降碳、减污、扩绿、增长，推进生态优先、节约集约、绿色低碳发展。
\begin{enumerate}[label=\Alph*., leftmargin=2em, itemsep=0.1em, topsep=0.15em]
\item 美丽中国建设
\item 生态环境整治
\item 乡村振兴
\item 环境保护
\end{enumerate}
\textbf{【答案】A}
\item 党的二十大报告指出，如期实现建军一百年奋斗目标，加快把人民军队建成世界一流军队，是全面建设社会主义现代化国家的\allowbreak\underline{\hspace{4em}}。
\begin{enumerate}[label=\Alph*., leftmargin=2em, itemsep=0.1em, topsep=0.15em]
\item 方针政策
\item 战略要求
\item 使命
\item 安全保障
\end{enumerate}
\textbf{【答案】B}
\item 党的二十大报告指出，全面加强人民军队党的建设，确保枪杆子永远听党指挥。健全贯彻\allowbreak\underline{\hspace{4em}}。
\begin{enumerate}[label=\Alph*., leftmargin=2em, itemsep=0.1em, topsep=0.15em]
\item 国家军委负责制
\item 军委主席负责制体制机制
\item 国防教育机制
\item 人民军队的组织体系
\end{enumerate}
\textbf{【答案】B}
\item 党的二十大报告指出，\allowbreak\underline{\hspace{4em}}是中国特色社会主义的伟大创举，是香港、澳门回归后保持长期繁荣稳定的最佳制度安排，必须长期坚持。
\begin{enumerate}[label=\Alph*., leftmargin=2em, itemsep=0.1em, topsep=0.15em]
\item 港人治港、澳人治澳
\item 特别行政区制度
\item “一国两制”
\item 中央全面管治权
\end{enumerate}
\textbf{【答案】C}
\item 党的二十大报告指出，全面准确、坚定不移贯彻“一国两制”、“港人治港”、“澳人治澳”、\allowbreak\underline{\hspace{4em}}的方针，坚持依法治港治澳，维护宪法和基本法确定的特别行政区宪制秩序。
\begin{enumerate}[label=\Alph*., leftmargin=2em, itemsep=0.1em, topsep=0.15em]
\item 高度自治
\item 全面自治
\item 特首领导
\item 中央全面管治
\end{enumerate}
\textbf{【答案】A}
\item 党的二十大报告指出，解决台湾问题、实现祖国完全统一，是党矢志不渝的历史任务，是全体中华儿女的共同愿望，是实现\allowbreak\underline{\hspace{4em}}的必然要求。
\begin{enumerate}[label=\Alph*., leftmargin=2em, itemsep=0.1em, topsep=0.15em]
\item 社会主义现代化强国
\item “一国两制”
\item 两岸繁荣发展
\item 中华民族伟大复兴
\end{enumerate}
\textbf{【答案】D}
\item 党的二十大报告指出，\allowbreak\underline{\hspace{4em}}是实现两岸统一的最佳方式，对两岸同胞和中华民族最有利。
\begin{enumerate}[label=\Alph*., leftmargin=2em, itemsep=0.1em, topsep=0.15em]
\item “和平统一、一国两制” 方针
\item “九二共识”
\item 两岸关系和平发展
\item 团结广大台湾同胞
\end{enumerate}
\textbf{【答案】A}
\item 党的二十大报告指出，中国始终坚持维护世界和平、促进共同发展的外交政策宗旨，致力于推动\allowbreak\underline{\hspace{4em}}。
\begin{enumerate}[label=\Alph*., leftmargin=2em, itemsep=0.1em, topsep=0.15em]
\item 全球经济健康发展
\item 构建人类命运共同体
\item 对外开放交流
\item 经济全球化
\end{enumerate}
\textbf{【答案】B}
\item 党的二十大报告指出，中国坚定奉行\allowbreak\underline{\hspace{4em}}的和平外交政策，始终根据事情本身的是非曲直决定自己的立场和政策，维护国际关系基本准则，维护国际公平正义。
\begin{enumerate}[label=\Alph*., leftmargin=2em, itemsep=0.1em, topsep=0.15em]
\item 不结盟
\item 独立自主
\item 平等
\item 不称霸、不扩张
\end{enumerate}
\textbf{【答案】B}
\item 党的二十大报告指出，中国坚持在\allowbreak\underline{\hspace{4em}}基础上同各国发展友好合作，推动构建新型国际关系，深化拓展平等、开放、合作的全球伙伴关系，致力于扩大同各国利益的汇合点。
\begin{enumerate}[label=\Alph*., leftmargin=2em, itemsep=0.1em, topsep=0.15em]
\item 和平共处五项原则
\item 经济互惠互利
\item 平等共处
\item 与邻为善、以邻为伴
\end{enumerate}
\textbf{【答案】A}
\item 党的二十大报告指出，中国坚持\allowbreak\underline{\hspace{4em}}的基本国策，坚定奉行互利共赢的开放战略，不断以中国新发展为世界提供新机遇，推动建设开放型世界经济，更好惠及各国人民。
\begin{enumerate}[label=\Alph*., leftmargin=2em, itemsep=0.1em, topsep=0.15em]
\item 对外开放
\item 和平共处五项原则
\item 反对保护主义
\item 不结盟、不称霸、不扩张
\end{enumerate}
\textbf{【答案】A}
\item 党的二十大报告指出，中国积极参与全球治理体系改革和建设，践行\allowbreak\underline{\hspace{4em}}的全球治理观，坚持真正的多边主义，推进国际关系民主化，推动全球治理朝着更加公正合理的方向发展。
\begin{enumerate}[label=\Alph*., leftmargin=2em, itemsep=0.1em, topsep=0.15em]
\item 平等互利
\item 互惠共赢
\item 平等交流、共同发展
\item 共商共建共享
\end{enumerate}
\textbf{【答案】D}
\item 党的二十大报告指出，\allowbreak\underline{\hspace{4em}}是世界各国人民前途所在。
\begin{enumerate}[label=\Alph*., leftmargin=2em, itemsep=0.1em, topsep=0.15em]
\item 共同繁荣、共同进步
\item 构建人类命运共同体
\item 平等交流、共同发展
\item 和平发展
\end{enumerate}
\textbf{【答案】B}
\item 党的二十大报告指出，全党必须牢记，\allowbreak\underline{\hspace{4em}}永远在路上，党的自我革命永远在路上，决不能有松劲歇脚、疲劳厌战的情绪，必须持之以恒推进全面从严治党，深入推进新时代党的建设新的伟大工程，以党的自我革命引领社会革命。
\begin{enumerate}[label=\Alph*., leftmargin=2em, itemsep=0.1em, topsep=0.15em]
\item 反腐
\item 全面从严治党
\item 改革开放
\item 全面建设社会主义现代化国家
\end{enumerate}
\textbf{【答案】B}
\item 党的二十大报告指出，坚持制度治党、依规治党，以党章为根本，以\allowbreak\underline{\hspace{4em}}为核心，完善党内法规制度体系，增强党内法规权威性和执行力，形成坚持真理、修正错误，发现问题、纠正偏差的机制。
\begin{enumerate}[label=\Alph*., leftmargin=2em, itemsep=0.1em, topsep=0.15em]
\item 依规问责
\item 巡视整改
\item 民主集中制
\item 从严治党
\end{enumerate}
\textbf{【答案】C}
\item 中国共产党第二十次全国代表大会，是在全党全国各族人民迈上\allowbreak\underline{\hspace{4em}}、向第二个百年奋斗目标进军的关键时刻召开的一次十分重要的大会。
\begin{enumerate}[label=\Alph*., leftmargin=2em, itemsep=0.1em, topsep=0.15em]
\item 全面建设社会主义现代化国家新征程
\item 全面建设中国特色社会主义现代化国家新征程
\item 全面建成小康社会
\item 建设中国特色社会主义现代化国家新征程
\end{enumerate}
\textbf{【答案】A}
\item 中国共产党第二十次全国代表大会的主题是：高举中国特色社会主义伟大旗帜，全面贯彻新时代中国特色社会主义思想，弘扬伟大建党精神，\allowbreak\underline{\hspace{4em}}，踔厉奋发、勇毅前行，为全面建设社会主义现代化国家、\allowbreak\underline{\hspace{4em}}而团结奋斗。
\begin{enumerate}[label=\Alph*., leftmargin=2em, itemsep=0.1em, topsep=0.15em]
\item 不忘初心、牢记使命；实现第二个百年奋斗目标
\item 自信自强、守正创新；全面推进中华民族伟大复兴
\item 不忘初心、牢记使命；全面推进中华民族伟大复兴
\item 自信自强、守正创新；实现第二个百年奋斗目标
\end{enumerate}
\textbf{【答案】B}
\item 党的二十大报告指出，中国共产党立志于中华民族千秋伟业，致力于人类和平与发展崇高事业，责任无比重大，使命无上光荣。全党同志务必\allowbreak\underline{\hspace{4em}}，务必\allowbreak\underline{\hspace{4em}}，务必\allowbreak\underline{\hspace{4em}}，坚定历史自信，增强历史主动，谱写新时代中国特色社会主义更加绚丽的华章。\\ ①不忘初心、牢记使命 ②谦虚谨慎、艰苦奋斗\\ ③勤劳勇敢、自强不息 ④敢于斗争、善于斗争
\begin{enumerate}[label=\Alph*., leftmargin=2em, itemsep=0.1em, topsep=0.15em]
\item ①②③
\item ①②④
\item ①③④
\item ②③④
\end{enumerate}
\textbf{【答案】B}
\item 党的二十大报告指出，我们隆重庆祝中国共产党成立一百周年、中华人民共和国成立七十周年，制定\allowbreak\underline{\hspace{4em}}历史决议，在全党开展党史学习教育，建成中国共产党历史展览馆，号召全党学习和践行伟大建党精神，在新的征程上更加坚定、更加自觉地牢记初心使命、开创美好未来。
\begin{enumerate}[label=\Alph*., leftmargin=2em, itemsep=0.1em, topsep=0.15em]
\item 第一个
\item 第二个
\item 第三个
\item 第四个
\end{enumerate}
\textbf{【答案】C}
\item 党的二十大报告指出，面对香港局势动荡变化，我们依照宪法和基本法有效实施对特别行政区的\allowbreak\underline{\hspace{4em}}，制定实施香港特别行政区维护国家安全法，落实“爱国者治港”原则，香港局势实现由乱到治的重大转折，深入推进粤港澳大湾区建设，支持香港、澳门发展经济、改善民生、保持稳定。
\begin{enumerate}[label=\Alph*., leftmargin=2em, itemsep=0.1em, topsep=0.15em]
\item 单方管治权
\item 全面管制权
\item 全面管理权
\item 全面管治权
\end{enumerate}
\textbf{【答案】D}
\item 中国特色社会主义最本质的特征是\allowbreak\underline{\hspace{4em}}。
\begin{enumerate}[label=\Alph*., leftmargin=2em, itemsep=0.1em, topsep=0.15em]
\item 突出政治建设
\item 坚持以人为本
\item 中国共产党领导
\item 集中统一领导
\end{enumerate}
\textbf{【答案】C}
\item 党的二十大报告指出，我们坚持精准扶贫、尽锐出战，打赢了人类历史上规模最大的脱贫攻坚战，全国八百三十二个贫困县全部摘帽，近\allowbreak\underline{\hspace{4em}}农村贫困人口实现脱贫，九百六十多万贫困人口实现易地搬迁，历史性地解决了绝对贫困问题，为全球减贫事业作出了重大贡献。
\begin{enumerate}[label=\Alph*., leftmargin=2em, itemsep=0.1em, topsep=0.15em]
\item 八千万
\item 九千万
\item 一亿
\item 一亿三千万
\end{enumerate}
\textbf{【答案】C}
\item 党的二十大报告指出，基础研究和原始创新不断加强，一些关键核心技术实现突破，战略性新兴产业发展壮大，载人航天、探月探火、深海深地探测、超级计算机、卫星导航、量子信息、核电技术、新能源技术、大飞机制造、生物医药等取得重大成果，进入\allowbreak\underline{\hspace{4em}}国家行列。
\begin{enumerate}[label=\Alph*., leftmargin=2em, itemsep=0.1em, topsep=0.15em]
\item 发达
\item 发展中
\item 创新型
\item 科技型
\end{enumerate}
\textbf{【答案】C}
\item 党的二十大报告指出，实践告诉我们，中国共产党为什么能，中国特色社会主义为什么好，归根到底是马克思主义行，是\allowbreak\underline{\hspace{4em}}的马克思主义行。拥有马克思主义科学理论指导是中国共产党坚定信仰信念、把握历史主动的根本所在。
\begin{enumerate}[label=\Alph*., leftmargin=2em, itemsep=0.1em, topsep=0.15em]
\item 中国化时代化
\item 时代化理论化
\item 理论化现代化
\item 中国化现代化
\end{enumerate}
\textbf{【答案】A}
\item 党的二十大报告指出，继续推进实践基础上的理论创新，首先要把握好新时代中国特色社会主义思想的\allowbreak\underline{\hspace{4em}}，坚持好、运用好贯穿其中的\allowbreak\underline{\hspace{4em}}。
\begin{enumerate}[label=\Alph*., leftmargin=2em, itemsep=0.1em, topsep=0.15em]
\item 世界观和方法论；立场观点方法
\item 世界观；立场观点方法
\item 世界观和方法论；立场观点
\item 世界观；立场观点
\end{enumerate}
\textbf{【答案】A}
\item \allowbreak\underline{\hspace{4em}}是全面建设社会主义现代化国家的首要任务。\allowbreak\underline{\hspace{4em}}是党执政兴国的第一要务。
\begin{enumerate}[label=\Alph*., leftmargin=2em, itemsep=0.1em, topsep=0.15em]
\item 高质量发展改革
\item 可持续发展发展
\item 可持续发展改革
\item 高质量发展发展
\end{enumerate}
\textbf{【答案】D}
\item 党的二十大报告指出，全方位夯实粮食安全根基，全面落实粮食安全党政同责，牢牢守住\allowbreak\underline{\hspace{4em}}亩耕地红线，逐步把永久基本农田全部建成高标准农田，深入实施种业振兴行动，强化农业科技和装备支撑，健全种粮农民收益保障机制和主产区利益补偿机制，确保中国人的饭碗牢牢端在自己手中。
\begin{enumerate}[label=\Alph*., leftmargin=2em, itemsep=0.1em, topsep=0.15em]
\item 十六亿
\item 十七亿
\item 十八亿
\item 十九亿
\end{enumerate}
\textbf{【答案】C}
\item “四项基本原则”不包括\allowbreak\underline{\hspace{4em}}。
\begin{enumerate}[label=\Alph*., leftmargin=2em, itemsep=0.1em, topsep=0.15em]
\item 坚持社会主义道路
\item 坚持人民民主专政
\item 坚持改革开放
\item 坚持马克思列宁主义、毛泽东思想
\end{enumerate}
\textbf{【答案】C}
\item 党的二十大报告指出，必须坚持\allowbreak\underline{\hspace{4em}}是第一生产力、\allowbreak\underline{\hspace{4em}}是第一资源、\allowbreak\underline{\hspace{4em}}是第一动力，深入实施科教兴国战略、人才强国战略、创新驱动发展战略，开辟发展新领域新赛道，不断塑造发展新动能新优势。
\begin{enumerate}[label=\Alph*., leftmargin=2em, itemsep=0.1em, topsep=0.15em]
\item 人才科技创新
\item 科技创新人才
\item 创新科技人才
\item 科技人才创新
\end{enumerate}
\textbf{【答案】D}
\item 党的领导主要是在\allowbreak\underline{\hspace{4em}}方面的领导。
\begin{enumerate}[label=\Alph*., leftmargin=2em, itemsep=0.1em, topsep=0.15em]
\item 政治、意识和组织
\item 政治、思想和组织
\item 政治、意识和发展
\item 政治、思想和发展
\end{enumerate}
\textbf{【答案】B}
\item 党的二十大报告指出，坚持\allowbreak\underline{\hspace{4em}}在我国现代化建设全局中的核心地位。
\begin{enumerate}[label=\Alph*., leftmargin=2em, itemsep=0.1em, topsep=0.15em]
\item 科技进步
\item 经济发展
\item 创新
\item 发展
\end{enumerate}
\textbf{【答案】C}
\item 党的二十大报告指出，全党要坚持全心全意为人民服务的根本宗旨，树牢群众观点，贯彻群众路线，尊重人民首创精神，坚持\allowbreak\underline{\hspace{4em}}，从群众中来、到群众中去，始终保持同人民群众的血肉联系，始终接受人民批评和监督，始终同人民同呼吸、共命运、心连心，不断巩固全国各族人民大团结，加强海内外中华儿女大团结，形成同心共圆中国梦的强大合力。
\begin{enumerate}[label=\Alph*., leftmargin=2em, itemsep=0.1em, topsep=0.15em]
\item 一切为了人民、一切依靠人民
\item 以人民利益为中心
\item 团结各族人民
\item 接受人民群众监督
\end{enumerate}
\textbf{【答案】A}
\item “起来，饥寒交迫的奴隶；起来，全世界受苦的人!满腔的热血已经沸腾，要为真理而斗争!”这是传遍欧洲、响彻世界的无产阶级的战歌——《国际歌》，歌词由欧仁·鲍狄埃于1871年6月1日创作，皮埃尔·狄盖特于1888年6月为其谱曲。我们要了解《国际歌》的创作背景，应该知道的历史是\allowbreak\underline{\hspace{4em}}。
\begin{enumerate}[label=\Alph*., leftmargin=2em, itemsep=0.1em, topsep=0.15em]
\item 法国巴黎公社
\item 法国里昂工人起义
\item 《共产党宣言》的发表
\item 十月革命
\end{enumerate}
\textbf{【答案】A}
\item 科技计划是根据国家科技发展规划和战略安排，以财政支持或以宏观政策调控、引导，由\allowbreak\underline{\hspace{4em}}组织和实施的科学研究或试验发展活动及相关的其他科学技术活动。《中共中央关于进一步全面深化改革、推进中国式现代化的决定》提出：“改进科技计划管理，强化\allowbreak\underline{\hspace{4em}}前瞻性、引领性布局。”
\begin{enumerate}[label=\Alph*., leftmargin=2em, itemsep=0.1em, topsep=0.15em]
\item 科研院所交叉前沿领域、重点领域
\item 科研院所基础研究领域、交叉前沿领域、重点领域
\item 政府部门交叉前沿领域、重点领域
\item 政府部门基础研究领域、交叉前沿领域、重点领域
\end{enumerate}
\textbf{【答案】D}
\item 中共中央、国务院印发《乡村全面振兴规划（2024-2027年）》。规划提出，到\allowbreak\underline{\hspace{4em}}年，乡村全面振兴取得决定性进展，农业现代化\allowbreak\underline{\hspace{4em}}，农村基本具备现代生活条件。
\begin{enumerate}[label=\Alph*., leftmargin=2em, itemsep=0.1em, topsep=0.15em]
\item 2027 基本实现
\item 2027 全面实现
\item 2035 基本实现
\item 2035 全面实现
\end{enumerate}
\textbf{【答案】C}
\item 金融是国民经济的\allowbreak\underline{\hspace{4em}}，与人民群众利益密切联系，关系中国式现代化建设全局。\allowbreak\underline{\hspace{4em}}是金融安全网第一道防线。
\begin{enumerate}[label=\Alph*., leftmargin=2em, itemsep=0.1em, topsep=0.15em]
\item 根基金融监管
\item 血脉金融监管
\item 根基金融机构公司治理
\item 血脉金融机构公司治理
\end{enumerate}
\textbf{【答案】B}
\item 《中共中央关于进一步全面深化改革、推进中国式现代化的决定》提出：“完善发挥\allowbreak\underline{\hspace{4em}}凝聚人心、汇聚力量政治作用的政策举措。”
\begin{enumerate}[label=\Alph*., leftmargin=2em, itemsep=0.1em, topsep=0.15em]
\item 党的领导
\item 群众路线
\item 统一战线
\item 独立自主
\end{enumerate}
\textbf{【答案】C}
\item 《中国共产党章程》是中国共产党的\allowbreak\underline{\hspace{4em}}的集中体现，是近百年党的建设宝贵经验的制度结晶，是贯穿党的一切活动和规范党员思想行为的最高准则。\\ ①统一意志 ②路线纲领 ③奋斗目标 ④指导思想
\begin{enumerate}[label=\Alph*., leftmargin=2em, itemsep=0.1em, topsep=0.15em]
\item ①②
\item ①②④
\item ①②③④
\item ①③④
\end{enumerate}
\textbf{【答案】C}
\item 党的二十大报告指出，育人的根本在于\allowbreak\underline{\hspace{4em}}。全面贯彻党的教育方针，落实\allowbreak\underline{\hspace{4em}}根本任务，培养德智体美劳全面发展的社会主义建设者和接班人。
\begin{enumerate}[label=\Alph*., leftmargin=2em, itemsep=0.1em, topsep=0.15em]
\item 立德；立德树人
\item 立德；立德育人
\item 文化教育；立德树人
\item 文化教育；立德育人
\end{enumerate}
\textbf{【答案】A}
\item 以下关于中共党史中重要会议说法错误的是\allowbreak\underline{\hspace{4em}}。
\begin{enumerate}[label=\Alph*., leftmargin=2em, itemsep=0.1em, topsep=0.15em]
\item 八七会议上否定了“陈独秀右倾错误”
\item 遵义会议结束了王明“左”倾冒险主义在中共中央的统治
\item 瓦窑堡会议确定建立抗日民族统一战线的方针政策
\item 十一届三中全会确立邓小平理论作为党的指导思想
\end{enumerate}
\textbf{【答案】D}
\item 有关遵义会议的表述不正确的是\allowbreak\underline{\hspace{4em}}。
\begin{enumerate}[label=\Alph*., leftmargin=2em, itemsep=0.1em, topsep=0.15em]
\item 取消了博古的军事指挥权
\item 确定了开展土地革命的方针
\item 确立了以毛泽东为核心的党中央的正确领导
\item 是中国共产党历史上生死攸关的转折点
\end{enumerate}
\textbf{【答案】B}
\item 现在，我们\allowbreak\underline{\hspace{4em}}的问题已经得到基本解决，但\allowbreak\underline{\hspace{4em}}的问题还没有得到根本解决，中国在世界上的形象很大程度上仍是\allowbreak\underline{\hspace{4em}}。
\begin{enumerate}[label=\Alph*., leftmargin=2em, itemsep=0.1em, topsep=0.15em]
\item “挨打”、“挨饿”　“挨骂”　“他塑”而非“自塑”
\item “挨打”、“挨骂”　“挨饿”　“他塑”而非“自塑”
\item “挨饿”、“挨骂”　“挨打”　“自塑”而非“他塑”
\item “挨打”、“挨饿”　“挨骂”　“自塑”而非“他塑”
\end{enumerate}
\textbf{【答案】A}
\item 在中华民族面临亡国灭种的至暗时刻，毛泽东同志的\allowbreak\underline{\hspace{4em}}以马克思主义者的深邃洞察力，科学分析敌我双方优势短板，准确预见战争走向进程，指明正确战略方针，有力驳斥了“亡国论”和“速胜论”，极大增强了国人夺取抗战胜利的决心和信心。这部彪炳史册的光辉著作，是中国共产党领导抗日战争的纲领性文献，对夺取抗战胜利具有重大指导意义。
\begin{enumerate}[label=\Alph*., leftmargin=2em, itemsep=0.1em, topsep=0.15em]
\item 《论持久战》
\item 《论联合政府》
\item 《中国革命和中国共产党》
\item 《星星之火，可以燎原》
\end{enumerate}
\textbf{【答案】A}
\item 党的二十届四中全会审议通过的核心文件是\allowbreak\underline{\hspace{4em}}。
\begin{enumerate}[label=\Alph*., leftmargin=2em, itemsep=0.1em, topsep=0.15em]
\item 《中共中央关于全面深化改革若干重大问题的决定》
\item 《中共中央关于进一步全面深化改革、推进中国式现代化的决定》
\item 《中共中央关于制定国民经济和社会发展第十五个五年规划的建议》
\item 《中共中央关于制定国民经济和社会发展第十四个五年规划和二〇三五年远景目标的建议》
\end{enumerate}
\textbf{【答案】C}
\item 2025年9月3日，纪念中国人民抗日战争暨世界反法西斯战争胜利80周年招待会在京隆重举行，习近平发表重要讲话。习近平指出，中国式现代化是走和平发展道路的现代化，中国永远是世界的\allowbreak\underline{\hspace{4em}}。
\begin{enumerate}[label=\Alph*., leftmargin=2em, itemsep=0.1em, topsep=0.15em]
\item 领导力量、团结力量、革新力量
\item 正义力量、和谐力量、复兴力量
\item 和平力量、稳定力量、进步力量
\item 核心力量、先锋力量、榜样力量
\end{enumerate}
\textbf{【答案】C}
\item 2025年中央经济工作会议明确指出，当前和今后一个时期经济发展的“首要任务”是\allowbreak\underline{\hspace{4em}}。
\begin{enumerate}[label=\Alph*., leftmargin=2em, itemsep=0.1em, topsep=0.15em]
\item 扩大内需
\item 推动高质量发展
\item 防范化解金融风险
\item 促进共同富裕
\end{enumerate}
\textbf{【答案】B}
\item 据新华社2026年9月1日报道，近日，中共中央宣传部等12部门联合印发通知，要求在9月集中组织开展“全民国防教育月”活动，主题为“\allowbreak\underline{\hspace{4em}}”。
\begin{enumerate}[label=\Alph*., leftmargin=2em, itemsep=0.1em, topsep=0.15em]
\item 弘扬伟大长征精神 同心共筑强大国防
\item 弘扬抗战精神 共筑国防长城
\item 依法开展国防教育 提升全民国防素养
\item 踔厉奋发强国防 勇毅前行向复兴
\end{enumerate}
\textbf{【答案】A}
\item “十五五”时期在基本实现社会主义现代化进程中具有承前启后的重要地位。实现社会主义现代化是一个\allowbreak\underline{\hspace{4em}}的历史过程，需要不懈努力、接续奋斗。
\begin{enumerate}[label=\Alph*., leftmargin=2em, itemsep=0.1em, topsep=0.15em]
\item 突破性探索，坚持守正创新
\item 持续性探索，坚持自我审视
\item 螺旋式递进、不断革新进步
\item 阶梯式递进、不断发展进步
\end{enumerate}
\textbf{【答案】D}
\end{enumerate}

\subsection*{二、填空题}

\begin{enumerate}[label=\textbf{\arabic*.}, leftmargin=2.5em, itemsep=0.5em]
\item 2022年10月16日，中国共产党第二十次全国代表大会在北京开幕。大会的主题是：高举中国特色社会主义伟大旗帜，全面贯彻新时代中国特色社会主义思想，弘扬伟大建党精神，自信自强、守正创新，踔厉奋发、勇毅前行，为\CJKunderline{全面建设社会主义现代化国家}、\CJKunderline{全面推进中华民族伟大复兴}而团结奋斗。
\item 经过不懈努力，党找到了\CJKunderline{自我革命}这一跳出治乱兴衰历史周期率的第二个答案，确保党永远不变质、不变色、不变味。
\item 中国共产党已走过百年奋斗历程。中国共产党立志于中华民族千秋伟业，致力于人类和平与发展崇高事业，责任无比重大，使命无上光荣。全党同志务必\CJKunderline{不忘初心、牢记使命}，务必\CJKunderline{谦虚谨慎、艰苦奋斗}，务必\CJKunderline{敢于斗争、善于斗争}，坚定历史自信，增强历史主动，谱写新时代中国特色社会主义更加绚丽的华章。
\item 中国式现代化，是中国共产党领导的社会主义现代化，既有各国现代化的共同特征，更有基于自己国情的中国特色。中国式现代化是\CJKunderline{人口规模巨大}的现代化，是\CJKunderline{全体人民共同富裕}的现代化，是\CJKunderline{物质文明和精神文明相协调}的现代化，是\CJKunderline{人与自然和谐共生}的现代化，是\CJKunderline{走和平发展道路的现代化}。
\item \CJKunderline{国家安全}是民族复兴的根基，\CJKunderline{社会稳定}是国家强盛的前提。
\item 全党必须牢记，\CJKunderline{坚持党的全面领导}是坚持和发展中国特色社会主义的必由之路，\CJKunderline{中国特色社会主义}是实现中华民族伟大复兴的必由之路，\CJKunderline{团结奋斗}是中国人民创造历史伟业的必由之路，\CJKunderline{贯彻新发展理念}是新时代我国发展壮大的必由之路，\CJKunderline{全面从严治党}是党永葆生机活力、走好新的赶考之路的必由之路。
\item \CJKunderline{和平与发展}是当代世界两大主题，必须坚持\CJKunderline{独立自主}的和平外交政策，为我国现代化建设争取有利的国际环境。
\item 全面建成社会主义现代化强国，总的战略安排是分两步走：从二〇二〇年到二〇三五年\CJKunderline{基本实现社会主义现代化}；从二〇三五年到本世纪中叶\CJKunderline{把我国建成富强民主文明和谐美丽的社会主义现代化强国}。未来五年是全面建设社会主义现代化国家开局起步的\CJKunderline{关键时期}。
\item \CJKunderline{全面依法治国}是国家治理的一场深刻革命，关系党执政兴国，关系人民幸福安康，关系党和国家长治久安。
\item 我们要坚持以\CJKunderline{人民安全}为宗旨、以\CJKunderline{政治安全}为根本、以\CJKunderline{经济安全}为基础、以\CJKunderline{军事科技文化社会安全}为保障、以\CJKunderline{促进国际安全}为依托，统筹外部安全和内部安全、国土安全和国民安全、传统安全和非传统安全、自身安全和共同安全，统筹维护和塑造国家安全，夯实国家安全和社会稳定基层基础，完善参与全球安全治理机制，建设更高水平的平安中国，以新安全格局保障新发展格局。
\item \CJKunderline{尊重自然}、\CJKunderline{顺应自然}、\CJKunderline{保护自然}，是全面建设社会主义现代化国家的内在要求。
\item \CJKunderline{如期实现建军一百年奋斗目标，加快把人民军队建成世界一流军队}，是全面建设社会主义现代化国家的战略要求。
\item \CJKunderline{“一国两制”}是中国特色社会主义的伟大创举，是香港、澳门回归后保持长期繁荣稳定的最佳制度安排，必须长期坚持。
\item \CJKunderline{人民民主}是社会主义的生命，是全面建设社会主义现代化国家的应有之义。\CJKunderline{全过程人民民主}是社会主义民主政治的本质属性，是最广泛、最真实、最管用的民主。
\item \CJKunderline{高质量发展}是全面建设社会主义现代化国家的首要任务。
\item \CJKunderline{教育}、\CJKunderline{科技}、\CJKunderline{人才}是全面建设社会主义现代化国家的基础性、战略性支撑。
\item 坚持大抓基层的鲜明导向，把\CJKunderline{基层党组织建设}成为有效实现党的领导的坚强战斗堡垒，激励党员发挥先锋模范作用，保持党员队伍先进性和纯洁性。
\item 六个“必须坚持”：必须坚持\CJKunderline{人民至上}，必须坚持\CJKunderline{自信自立}，必须坚持\CJKunderline{守正创新}，必须坚持\CJKunderline{问题导向}，必须坚持\CJKunderline{系统观念}，必须坚持\CJKunderline{胸怀天下}。
\item 1848年发表的《\CJKunderline{共产党宣言}》是马克思主义诞生的标志。
\item 马克思主义主要包括\CJKunderline{马克思主义哲学}、\CJKunderline{马克思主义政治经济学}和\CJKunderline{科学社会主义}三个组成部分。
\item \CJKunderline{先进性}和\CJKunderline{纯洁性}是马克思主义政党的本质属性。
\item \CJKunderline{习近平新时代中国特色社会主义思想}是马克思主义中国化最新成果，是党和人民实践经验和集体智慧的结晶。
\item 1919年，\CJKunderline{李大钊}在\CJKunderline{新青年杂志}发表了\CJKunderline{《我的马克思主义观》}一文，是中国第一篇比较系统介绍马克思主义的著名论文。
\item 根据马克思主义理论，国家就其起源来说，它是\CJKunderline{阶级矛盾不可调和的产物}。
\item \CJKunderline{《中国共产党章程》}是最根本的党内法规，是管党治党的总规矩。
\item 《中国共产党党章》中规定，中国共产党是中国\CJKunderline{工人阶级}的先锋队，同时是\CJKunderline{中国人民}和\CJKunderline{中华民族}的先锋队。
\item 中国共产党党员是中国\CJKunderline{工人阶级}的有共产主义觉悟的先锋战士。中国共产党党员必须\CJKunderline{全心全意为人民服务}，不惜牺牲个人的一切，为实现共产主义奋斗终身。
\item 中国共产党党员永远是\CJKunderline{劳动人民}的普通一员。除了法律和政策规定范围内的个人利益和工作职权以外，所有共产党员都不得谋求任何私利和特权。
\item 党的最高理想和最终目标是\CJKunderline{实现共产主义}。
\item 党的最高领导机关，是党的\CJKunderline{全国代表大会}和它所产生的中央委员会。
\item 中国共产党的根本组织制度是\CJKunderline{民主集中制}。
\item 党的纪律主要包括\CJKunderline{政治纪律}、\CJKunderline{组织纪律}、\CJKunderline{廉洁纪律}、\CJKunderline{群众纪律}、\CJKunderline{工作纪律}、\CJKunderline{生活纪律}。
\item 党的纪律处分有五种：\CJKunderline{警告、严重警告、撤销党内职务、留党察看、开除党籍}。
\item 预备党员的权利，除了没有\CJKunderline{表决权}、\CJKunderline{选举权}和\CJKunderline{被选举权}以外，也同正式党员一样。党员如果没有正当理由，连续\CJKunderline{六}个月不参加党的组织生活，或不交纳党费，或不做党所分配的工作，就被认为是自行脱党。
\item \CJKunderline{人民立场}是中国共产党的根本政治立场。
\item 中国共产党的宗旨是\CJKunderline{全心全意为人民服务}。
\item 我党的最大政治优势是\CJKunderline{密切联系群众}。
\item 党的思想路线是\CJKunderline{一切从实际出发，理论联系实际，实事求是，在实践中检验真理和发展真理}。
\item “两学一做”是指\CJKunderline{学党章党规}、\CJKunderline{学系列讲话}，\CJKunderline{做合格党员}。
\item 党的领导主要是\CJKunderline{政治}、\CJKunderline{思想}和\CJKunderline{组织}领导。
\item 中国特色社会主义最本质的特征是\CJKunderline{中国共产党领导}。
\item 中国共产党人追求的\CJKunderline{共产主义}最高理想，只有在\CJKunderline{社会主义社会}充分发展和高度发达的基础上才能实现。\CJKunderline{社会主义制度}的发展和完善是一个长期的历史过程。
\item 新形势下党面临的四大考验分别是\CJKunderline{长期执政考验}、\CJKunderline{改革开放考验}、\CJKunderline{市场经济考验}、\CJKunderline{外部环境考验}。
\item \CJKunderline{党风问题、党同人民群众联系问题}是关系党生死存亡的问题。
\item 发展是中国共产党执政兴国的第一要务。必须坚持\CJKunderline{以人民为中心}的发展思想，坚持\CJKunderline{创新、协调、绿色、开放、共享}的发展理念。
\item 学“四史”，这里的“四史”是指：\CJKunderline{党史}、\CJKunderline{新中国史}、\CJKunderline{改革开放史}、\CJKunderline{社会主义发展史}。
\item 近代中国压在人民头上的三座大山是\CJKunderline{帝国主义}、\CJKunderline{封建主义}和\CJKunderline{官僚资本主义}。
\item 中国近代历史是一部耻辱史，也是悲壮的抗争史。为着救亡的呐喊，为着复兴的梦想，多少主义出场，又都破灭；多少道路探索，又都碰壁；多少政党登台，又都谢幕。这些悲壮的抗争和艰辛的求索证明\CJKunderline{资本主义道路在中国走不通，只有社会主义才能救中国}。
\item 中国新民主主义革命胜利的三大法宝：\CJKunderline{统一战线}、\CJKunderline{武装斗争}、\CJKunderline{党的建设}。
\item “工农武装割据”是指在中国共产党领导下，以\CJKunderline{武装斗争}为主要形式，以\CJKunderline{土地革命}为中心内容，以\CJKunderline{农村革命根据地}为战略阵地的三者密切结合。
\item \CJKunderline{瓦窑堡会议}是于1935年召开的中共中央政治局扩大会议，会议明确提出建立抗日民族统一战线的基本策略任务。此次会议对推动抗日民族统一战线和第二次国共合作的建立具有重要意义。
\item \CJKunderline{洛川会议}是1937年中国共产党在历史转折关头召开的一次重要会议。它制定了中国共产党的全面抗战路线，为中国共产党和全国人民指明了抗战的正确方向。
\item 毛泽东思想的活的灵魂三个基本方面：\CJKunderline{实事求是}、\CJKunderline{群众路线}、\CJKunderline{独立自主}。
\item 邓小平指出，社会主义的最大优越性是\CJKunderline{共同富裕}。
\item 以邓小平同志为核心的中国共产党第二代中央领导集体探索并初步回答了“什么是社会主义、怎样建设社会主义”的问题。1992年初，邓小平在“南方谈话” 中提出：“社会主义的本质，是\CJKunderline{解放生产力，发展生产力，消灭剥削，消除两极分化，最终达到共同富裕}。
\item 中华民族进行抗日战争抗战从\CJKunderline{1931}年九一八事变开始算起，至\CJKunderline{1945}年结束，共十四年。2014年，十二届全国人大常委会第七次会议经表决通过，将\CJKunderline{9月3日}确定为中国人民抗日战争胜利纪念日。
\item 第一次明确地把毛泽东思想确立为全党的指导思想，并写入党章的会议是\CJKunderline{中共七大}。
\item 七届二中全会首次提出加强执政党的建设，提出“两个务必”：务必使同志们继续地保持\CJKunderline{谦虚谨慎不骄不躁}的作风和务必使同志们继续地保持\CJKunderline{艰苦奋斗}的作风。
\item 到1952年底，广大新解放区的\CJKunderline{土地改革}基本完成，从根本上铲除了中国封建制度的根基。
\item \CJKunderline{1992}年，我国经济体制改革正式确立以建立社会主义市场经济体制的为目标。
\item \CJKunderline{2017}年，中国共产党第十九次全国代表大会在北京召开。会议主题是：不忘初心，牢记使命，高举\CJKunderline{中国特色社会主义}伟大旗帜，决胜全面建成小康社会，夺取新时代中国特色社会主义伟大胜利，为实现中华民族伟大复兴的中国梦不懈奋斗。
\item 从1949年10月中华人民共和国成立到1956年，中国共产党领导全国各族人民迅速恢复了国民经济并开展了有计划的经济建设，在全国绝大部分地区基本上完成了三大改造，确立了\CJKunderline{社会主义基本制度}。
\item 改革开放以来我们取得一切成绩和进步的根本原因，归结起来就是：开辟了\CJKunderline{中国特色社会主义道路}，\CJKunderline{形成了中国特色社会主义理论体系}，确立了\CJKunderline{中国特色社会主义制度}，发展了\CJKunderline{中国特色社会主义文化}。
\item 中国特色社会主义道路是\CJKunderline{实现途径}，中国特色社会主义理论体系是\CJKunderline{行动指南}，中国特色社会主义制度是\CJKunderline{根本保障}，中国特色社会主义文化是\CJKunderline{精神力量}，四者统一于中国特色社会主义伟大实践。
\item “三严三实”：\CJKunderline{严以修身}、\CJKunderline{严以用权}、\CJKunderline{严以律己}、\CJKunderline{谋事要实}、\CJKunderline{创业要实}、\CJKunderline{做人要实}。
\item 形势越变化、党和人民事业越发展，越要重视干部培养。培养干部，要抓好\CJKunderline{党性教育}这个核心，抓好\CJKunderline{道德建设}这个基础，加强宗旨意识、公仆意识教育。
\item 支部“三会一课”：党员大会、支部委员会、党小组会、党课
\item 提高党领导经济工作科学化水平。要围绕经济社会发展重大问题加强学习和调研，提高把握和运用市场经济规律、自然规律、\CJKunderline{社会发展规律}能力，坚持\CJKunderline{全球视野}，增强战略思维能力，做到厚积薄发。
\item 新中国自行研制的第一架轻型旅客机叫做\CJKunderline{北京一号}。1958年9月22日，由北京航空学院（现北京航空航天大学）师生研制的亚洲第一枚探空火箭成功发射，这枚火箭的名字是\CJKunderline{北京二号}。
\item 坚持把发展经济的着力点放在\CJKunderline{实体经济}上，推进新型工业化，加快建设制造强国、质量强国、航天强国、交通强国、网络强国、数字中国。
\item 北航学生党校行为规范：\CJKunderline{学习理论}，\CJKunderline{坚定信念}，\CJKunderline{牢记宗旨}，\CJKunderline{勇挑重担}，\CJKunderline{艰苦朴素}，\CJKunderline{勤奋好学}，\CJKunderline{严于律己}，\CJKunderline{争做模范}。
\item 北航校训：\CJKunderline{德才兼备}，\CJKunderline{知行合一}。
\item 2024年9月27日，习近平在全国民族团结进步表彰大会上指出，中国共产党历来高度重视民族工作。一百多年来，我们坚持把马克思主义民族理论同\CJKunderline{中国民族问题具体实际}相结合、同中华优秀传统文化相结合，创造性地走出了一条中国特色解决民族问题的正确道路。在这条道路上，党团结带领全国各族人民实现了\CJKunderline{民族独立和人民解放}，开创了发展各民族\CJKunderline{平等团结互助和谐}关系的新局面，推动民族地区经济社会发展和少数民族群众生活取得前所未有的进步，我国少数民族面貌、民族地区面貌、民族关系面貌、中华民族面貌发生了翻天覆地的历史性巨变。
\item 坚持党的文化领导权，必须把坚持和加强党对宣传思想文化工作的\CJKunderline{全面领导}落到实处。党的领导对于宣传思想文化工作起着\CJKunderline{定性定调定向}的根本性作用，坚持和加强党对宣传思想文化工作的全面领导是马克思主义政党的突出特点和\CJKunderline{优良传统}，也是做好宣传思想文化工作、推动新时代中国特色社会主义文化\CJKunderline{繁荣发展}的根本要求。
\item 2024年，党中央坚持以党的自我革命引领社会革命，一以贯之推进\CJKunderline{全面从严治党}，取得新的进展和成效。全党在自我革命这一重大问题上认识更加清醒、行动更加坚定；扎实开展党纪学习教育，全党同志遵规守纪意识明显增强；持续加大\CJKunderline{反腐败}力度，铲除腐败滋生的土壤和条件；深入开展群众身边不正之风和腐败问题集中整治，解决了一大批群众反映强烈的突出问题；深化\CJKunderline{制度治党}、\CJKunderline{依规治党}，进一步提高了党领导反腐败斗争、推进全面从严治党的能力。
\item 中共中央政治局2025年2月28日上午就建设更高水平平安中国进行第十九次集体学习。中共中央总书记习近平在主持学习时强调，要把捍卫国家\CJKunderline{政治安全}摆在首位，坚定维护国家\CJKunderline{政权安全、制度安全、意识形态安全}。
\item 2024年是实现\CJKunderline{“十四五”}规划目标任务的关键一年。面对外部压力加大、内部困难增多的复杂严峻形势，以习近平同志为核心的党中央团结带领全党全国各族人民，沉着应变、综合施策，经济运行总体平稳、稳中有进，高质量发展扎实推进，经济社会发展主要目标任务即将顺利完成。\CJKunderline{新质生产力}稳步发展，\CJKunderline{改革开放}持续深化，\CJKunderline{重点领域风险}化解有序有效，民生保障扎实有力，中国式现代化迈出新的坚实步伐。
\item 新质生产力是\CJKunderline{创新}起主导作用，摆脱传统经济增长方式、生产力发展路径，具有\CJKunderline{高科技}、\CJKunderline{高效能}、\CJKunderline{高质量}特征，符合新发展理念的先进生产力质态。它由技术革命性突破、生产要素创新性配置、产业深度转型升级而催生，以劳动者、劳动资料、劳动对象及其优化组合的跃升为基本内涵，以全要素生产率大幅提升为核心标志，特点是\CJKunderline{创新}，关键在\CJKunderline{质优}，本质是\CJKunderline{先进生产力}。
\item 在2015年纪念中国人民抗日战争暨世界反法西斯战争胜利70周年大会上，习近平总书记发表了“让我们共同铭记历史所启示的伟大真理：\CJKunderline{正义必胜}！\CJKunderline{和平必胜}！\CJKunderline{人民必胜}！”这一振聋发聩的宣言。十年后，在2025年的阅兵式上，这一宣言再次以嘹亮口号响彻于世界人民面前，以鲜明标语高悬于天安门上空，深刻揭示了历史发展的规律，道出了中国人民与世界人民的共同心声。
\item 习近平总书记在纪念中国人民抗日战争暨世界反法西斯战争胜利80周年大会上讲话提到：历史\CJKunderline{承载}过去，也\CJKunderline{启迪}未来。新时代新征程，全国各族人民要在中国共产党坚强领导下，坚持马克思列宁主义、毛泽东思想、邓小平理论、“三个代表”重要思想、科学发展观，全面贯彻新时代中国特色社会主义思想，坚定不移走\CJKunderline{中国特色社会主义道路}，传承和弘扬\CJKunderline{伟大抗战精神}，踔厉奋发、勇毅前行，为以中国式现代化全面推进强国建设、民族复兴伟业而团结奋斗！
\item 中国共产党第二十届中央委员会第\CJKunderline{四}次全体会议于2025年10月20日至23日在北京举行。全会审议通过了《中共中央关于制定国民经济和社会发展\CJKunderline{第十五个五年规划}的建议》。
\item 2026年是\CJKunderline{“十五五”}规划的\CJKunderline{开局之年}。我们要坚持\CJKunderline{稳中求进}工作总基调，实施更加积极有为的宏观政策，推动经济实现质的有效提升和量的合理增长。
\item “十五五”时期经济社会发展必须遵循以下原则：坚持党的全面领导，坚持人民至上，坚持高质量发展，坚持全面深化改革，坚持\CJKunderline{有效市场}和\CJKunderline{有为政府}相结合，坚持统筹发展和安全。
\item 中共中央党史和文献研究院编辑的《习近平文化文选》第一卷收入《提高国家文化软实力，展示中华文化独特魅力》（2013年12月30日）。文中指出：提高国家文化软实力，不仅关系我国在世界文化格局中的定位，而且关系我国国际地位和国际影响力，关系“两个一百年”奋斗目标和中华民族伟大复兴中国梦的实现。要努力夯实\CJKunderline{国家文化软实力}的根基、传播当代\CJKunderline{中国价值}观念、展示\CJKunderline{中华文化}独特魅力、提高\CJKunderline{国际话语权}，让当代中国形象在世界上不断树立和闪亮起来。
\item “十五五”时期经济社会发展的主要目标：\CJKunderline{高质量发展}取得显著成效，\CJKunderline{科技自立自强水平}大幅提高，进一步全面深化改革取得新突破，社会文明程度明显提升，人民生活品质不断提高，美丽中国建设取得新的重大进展，\CJKunderline{国家安全屏障}更加巩固。在此基础上再奋斗五年，到二〇三五年实现我国经济实力、科技实力、国防实力、综合国力和国际影响力大幅跃升，人均国内生产总值达到中等发达国家水平，人民生活更加幸福美好，基本实现社会主义现代化。
\end{enumerate}

\subsection*{三、简答题}

\begin{enumerate}[label=\textbf{\arabic*.}, leftmargin=2.5em, itemsep=0.8em]
\item 请默写入党誓词。
\par \CJKunderline{我志愿加入中国共产党，拥护党的纲领，遵守党的章程，履行党员义务，执行党的决定，严守党的纪律，保守党的秘密，对党忠诚，积极工作，为共产主义奋斗终身，随时准备为党和人民牺牲一切，永不叛党。}
\item 《国际歌》第一段歌词。
\par \CJKunderline{起来，饥寒交迫的奴隶！起来，全世界受苦的人！满腔的热血已经沸腾，要为真理而斗争！旧世界打个落花流水，奴隶们起来，起来！不要说我们一无所有，我们要做天下的主人！这是最后的斗争，团结起来到明天，英特纳雄耐尔就一定要实现！这是最后的斗争，团结起来到明天，英特纳雄耐尔就一定要实现！}
\item 请写出“四个自信”和“两个维护”的内容。
\par \CJKunderline{“四个自信”：中国特色社会主义道路自信、理论自信、制度自信、文化自信；}
\par \CJKunderline{“两个维护”：坚决维护习近平总书记党中央的核心、全党的核心地位，坚决维护党中央权威和集中统一领导。}
\item 请写出“五位一体”总体布局和“四个全面”战略布局的内容。
\par \CJKunderline{“五位一体”：经济建设、政治建设、文化建设、社会建设、生态文明建设五位一体。}
\par \CJKunderline{“四个全面”：全面建设社会主义现代化国家、全面深化改革、全面依法治国、全面从严治党。}
\item 请写出新质生产力的核心标志、关键、本质、核心要素。
\par \CJKunderline{新质生产力以全要素生产率大幅提升为核心标志，关键在质优，本质是先进生产力。科技创新能够催生新产业、新模式、新动能，是发展新质生产力的核心要素。}
\item 请写出中国共产党的行动指南。
\par \CJKunderline{马克思列宁主义，毛泽东思想，邓小平理论，“三个代表”重要思想，科学发展观，习近平新时代中国特色社会主义思想。}
\item 请写出“四个意识”和“四个伟大”内容。
\par \CJKunderline{“四个意识”：政治意识、大局意识、核心意识、看齐意识；“四个伟大”：实现伟大梦想，进行伟大斗争，建设伟大工程，推进伟大事业。}
\item 请写出“四风”问题的内容。
\par \CJKunderline{形式主义、官僚主义、享乐主义、奢靡之风。}
\item 请写出党的三大作风和党的三大法宝。
\par \CJKunderline{党的三大作风：理论联系实际、密切联系群众、批评与自我批评；}
\par \CJKunderline{党的三大法宝：统一战线、武装斗争、党的建设。}
\item 请写出党的思想路线和群众路线的基本内容。
\par \CJKunderline{思想路线：一切从实际出发，理论联系实际，实事求是，在实践中检验真理和发展真理。}
\par \CJKunderline{群众路线：一切为了群众，一切依靠群众，从群众中来，到群众中去，把党的正确主张变为群众的自觉行动。}
\item 请写出和平共处五项原则。
\par \CJKunderline{互相尊重主权和领土完整、互不侵犯、互不干涉内政、平等互利、和平共处。}
\item 请写出党面临的四个危险。
\par \CJKunderline{精神懈怠危险、能力不足危险、脱离群众危险、消极腐败危险。}
\item 请说出二十大报告中中国式现代化的本质要求？
\par \CJKunderline{中国式现代化的本质要求是：坚持中国共产党领导，坚持中国特色社会主义，实现高质量发展，发展全过程人民民主，丰富人民精神世界，实现全体人民共同富裕，促进人与自然和谐共生，推动构建人类命运共同体，创造人类文明新形态。}
\item 请说出二十大报告中“十年来我们经历了对党和人民事业具有重大现实意义和深远历史意义的三件大事”？
\par \CJKunderline{一是迎来了中国共产党成立一百周年，二是中国特色社会主义进入新时代，三是完成脱贫攻坚、全面建成小康社会的历史任务，实现第一个百年奋斗目标。}
\item 请说出“党的中心任务”？
\par \CJKunderline{党的二十大强调，从现在起，中国共产党的中心任务就是团结带领全国各族人民全面建成社会主义现代化强国、实现第二个百年奋斗目标，以中国式现代化全面推进中华民族伟大复兴。}
\end{enumerate}

\begin{center}\large 希望同学们在新的学期里只争朝夕，不负韶华！祝愿每位同学都能以优异的成绩进入党校学习！\end{center}
\end{document}
